% =====================================================================
%  Submission version.
%  Compiles standalone with pdflatex; run twice for references.
% =====================================================================
\documentclass[11pt,a4paper]{article}

\usepackage{amsmath}
\usepackage{amssymb}
\usepackage{amsthm}
\usepackage[margin=1in]{geometry}
\usepackage[colorlinks=true,linkcolor=blue,citecolor=blue]{hyperref}

\theoremstyle{plain}
\newtheorem{theorem}{Theorem}
\newtheorem{corollary}[theorem]{Corollary}
\newtheorem{proposition}[theorem]{Proposition}
\newtheorem{lemma}[theorem]{Lemma}
\theoremstyle{definition}
\newtheorem{hypothesis}[theorem]{Hypothesis}
\theoremstyle{remark}
\newtheorem{remark}[theorem]{Remark}

\DeclareMathOperator{\dist}{dist}
\DeclareMathOperator{\cone}{cone}
\newcommand{\Pp}{P^{+}}
\newcommand{\Pm}{P^{-}}
\newcommand{\thr}{\mathrm{thr}}
\newcommand{\li}{\mathrm{li}}
\newcommand{\Li}{\mathrm{Li}}
\newcommand{\vB}{\mathcal{E}}
\newcommand{\vR}{\mathcal{O}}
\newcommand{\band}{\mathcal{B}}
\newcommand{\eps}{\varepsilon}

% typographic slack: let TeX stretch a line rather than run into the margin
\emergencystretch=3em

\title{The parity-separation constant of level-$Q$ divisor weights
tends to zero}
\author{Youngrok Lee\\
\small Independent researcher, Republic of Korea\\ \small ORCID \href{https://orcid.org/0009-0005-0409-5267}{0009-0005-0409-5267}}
\date{}

\begin{document}
\maketitle

\begin{abstract}
For an integer $N$ put $Q=\lfloor\sqrt N\rfloor$, let $q_{1}$ be the least prime exceeding $Q$, and let
$\band$ be the set of squarefree $Q$-friable integers in $(\lfloor N/q_{1}\rfloor,N]$. Let $\kappa(N)$ be
the largest average negativity, over the odd-$\omega$ part of $\band$, of a divisor weight of level $Q$
that is nonnegative on the even-$\omega$ part, normalised by the $\ell^{2}$ norm of its coefficients and by
the $\ell^{2}$ norm of the parity defect of the band; equivalently
$\kappa(N)=\dist_{2}(c_{\vR},K_{Q})/\|c_{\vB}-c_{\vR}\|_{2}$, where $c_{\vB}$ and $c_{\vR}$ are the
divisor-incidence count vectors of the two parities and $K_{Q}$ is the cone generated by the
divisor-incidence columns of the even-$\omega$ part. Weighting all even-$\omega$ columns equally gives
$\kappa(N)\le1$, so $\kappa(N)$ measures how much a nonnegative reweighting improves on that baseline. We
prove that
\[
  \kappa(N)\ \ll_{A}\ (\log Q)^{-A}\qquad\text{for every fixed }A>0;
\]
in particular $\kappa(N)\to0$. The quantitative form of this, with a parameter $\beta\ge4$ fixed before
$N$ grows, is a two-term bound whose leading coefficients are evaluated and whose threshold of validity
is not: at $\beta=6$ it reads $\kappa(N)\le(2.01317+\eps)/(\log Q\sqrt{\log\log Q})$ for all
sufficiently large $N$, the constant being recorded as an example rather than as an optimal one. The proof is
a direct positive construction: explicit nonnegative weights on the even-$\omega$ columns, grouped by the
small-prime part $a_{Y}(\ell)=\prod_{p\mid\ell,\,p\le Y}p$ and truncated at $a_{Y}(\ell)\le\sqrt Q$. The
one arithmetic input it needs is that, after fixing the small-prime part, the two $\omega$-parities of the
rough cofactor occur in bounded ratio; we prove this. The elementary observation that drives the signed
estimate is that $N<q_{1}^{2}$, so an integer $n\le N$ has at most one prime factor exceeding $Q$; the
classical large-prime removal is then exact with no upper cutoff on the cofactor, and the required signed
two-cutoff estimate follows from a quantitative prime number theorem together with the matching M\"obius
bound. The result is effective in principle, but the constant implied by $\ll_{A}$ is not evaluated. The
bound is an upper bound only: it says nothing about whether $\kappa(N)$ is positive, nor whether the
rate is optimal.
\end{abstract}

{\small\noindent 2020 Mathematics Subject Classification: 11N35 (primary); 11N36, 11N25,
90C05 (secondary). Keywords: parity problem, divisor weights, convex cones, linear programming
duality, certified computation.\par}

\section{Introduction}\label{sec:intro}

\subsection{The object}\label{sec:object}
Throughout, $N$ is an integer and
\begin{equation}\label{eq:param}
  Q=\lfloor\sqrt N\rfloor,\qquad L=\log Q,\qquad
  q_{1}=\min\{p\ \text{prime}:p>Q\},\qquad \thr=\lfloor N/q_{1}\rfloor,
\end{equation}
with $\log$ the natural logarithm. For $n>1$ let $\Pm(n)$, $\Pp(n)$ be the least and greatest prime factors
of $n$, with $\Pm(1)=\infty$ and $\Pp(1)=1$; $\omega(n)$ counts distinct prime factors, and $\mu$ is the
M\"obius function. The \emph{band} is
\begin{equation}\label{eq:band}
  \band=\bigl\{n\in(\thr,N]:\ \mu^{2}(n)=1,\ \Pp(n)\le Q\bigr\},
\end{equation}
split by the parity of $\omega$ into
\begin{equation}\label{eq:parts}
  \vB=\{n\in\band:\omega(n)\ \text{even}\},\qquad
  \vR=\{n\in\band:\omega(n)\ \text{odd}\},\qquad R=|\vR|.
\end{equation}
Index coordinates by the \emph{rows}, the squarefree $d\le Q$. With $c_{\vB}(d)=\#\{n\in\vB:d\mid n\}$ and
$c_{\vR}(d)=\#\{n\in\vR:d\mid n\}$, and assuming $R>0$, set
\begin{equation}\label{eq:tp}
  t_{d}=\frac{c_{\vR}(d)}R,\qquad
  p_{d}=\frac{c_{\vB}(d)-c_{\vR}(d)}R=\frac1R\sum_{n\in\band,\ d\mid n}\mu(n),
\end{equation}
the second equality because $\mu(n)=(-1)^{\omega(n)}$ on squarefree $n$. For $\ell\in\vB$ let
$A_{\ell}=(\mathbf1_{d\mid\ell})_{d}$ be its divisor-incidence column, let $A$ be the matrix with these
columns, put
\begin{equation}\label{eq:cone}
  K_{Q}=\cone\{A_{\ell}:\ell\in\vB\},
\end{equation}
and define
\begin{equation}\label{eq:kappa}
  \kappa(N)=\frac{\dist_{2}(t,K_{Q})}{\|p\|_{2}},
\end{equation}
both norms being Euclidean on the row space. The vector $p$ is the \emph{parity defect} of the band:
$p_{1}$ is the normalised M\"obius sum over $\band$, and $p_{d}$ its analogue along the row $d$.

\paragraph{The baseline, and an $R$-free form.}
Both $t=c_{\vR}/R$ and $p=(c_{\vB}-c_{\vR})/R$ carry the same factor $1/R$, where
$c_{\vB}=(c_{\vB}(d))_{d}$ and $c_{\vR}=(c_{\vR}(d))_{d}$, and $\dist_{2}(\cdot,K_{Q})$ is homogeneous
because $K_{Q}$ is a cone. So $R$ cancels in \eqref{eq:kappa}:
\begin{equation}\label{eq:Rfree}
  \kappa(N)=\frac{\dist_{2}(c_{\vR},K_{Q})}{\|c_{\vB}-c_{\vR}\|_{2}} .
\end{equation}
This form needs $R>0$ and, as in \eqref{eq:kappa}, a nonzero denominator: $c_{\vB}\neq c_{\vR}$, which is
the condition $p\neq0$ cleared of the factor $1/R$. What \eqref{eq:Rfree} removes is the normalisation by
$R$, not either of those two conditions. The quantitative bound $N/R=O(1)$ of
Proposition~\ref{prop:norm} is a separate input, used in the reduction of \S\ref{sec:reduction} and
nowhere in the definition. Taking the uniform weight $x^{(0)}_{\ell}=1/R$ on every even-$\omega$ column
gives $(Ax^{(0)})_{d}=c_{\vB}(d)/R=t_{d}+p_{d}$, so $\dist_{2}(t,K_{Q})\le\|p\|_{2}$ and
\begin{equation}\label{eq:kappale1}
  0\ \le\ \kappa(N)\ \le\ 1 .
\end{equation}
This fixes the scale and gives the quantity its reading: \emph{the denominator of \eqref{eq:Rfree} is the
error made by weighting all even-$\omega$ columns equally, the numerator is what the best nonnegative
reweighting achieves, and $\kappa(N)$ is the ratio.} What this paper proves is that the optimum beats the
uniform baseline by more than any fixed power of $\log Q$.

Two remarks on the definition. First, \eqref{eq:tp} and \eqref{eq:kappa} require $R>0$ and $p\neq0$;
Proposition~\ref{prop:norm} shows that both hold once $N$ is large, and every statement below is for such
$N$. Second, although the level $Q$ enters the name of the constant, \eqref{eq:band}--\eqref{eq:kappa} are
determined by $N$ and not by $Q$: the $2Q+1$ integers $N$ with $Q^{2}\le N<(Q+1)^{2}$ share the level $Q$
but need not share the band, although they may (at $N=10$ and $N=11$ alike, $Q=3$, $\thr=2$ and the band
is $\{3,6\}$).
For instance at $Q=5$ the odd part of the band is $\{5\}$ at $N=25$, whose row vector is not in the cone
generated by the columns of $6,10,15$, whereas at $N=30$ it is $\{5,30\}$ and
$t=\tfrac12(A_{10}+A_{15})\in K_{Q}$, so that $\kappa(25)>0=\kappa(30)$. We therefore write $\kappa(N)$.

\subsection{What \texorpdfstring{$\kappa$}{kappa} measures}\label{sec:measures}
Let $\lambda=(\lambda_{d})_{d\le Q,\ \mu^{2}(d)=1}$ be a weight of \emph{level} $Q$, let
$\theta_{\lambda}(n)=\sum_{d\mid n,\ d\le Q}\lambda_{d}$ be its divisor sum, and let $\mathbb E_{\vR}$
denote the mean over $\vR$.

\begin{proposition}\label{prop:dual}
Let $N$ be such that $R>0$. Then
\[
  \dist_{2}(t,K_{Q})=\max\bigl\{-\mathbb E_{\vR}\theta_{\lambda}:\
    \|\lambda\|_{2}\le1,\ \theta_{\lambda}(\ell)\ge0\ \text{for all }\ell\in\vB\bigr\}.
\]
If $t\in K_{Q}$ the maximum is $0$, attained at $\lambda=0$. If $t\notin K_{Q}$ it is positive and attained
at a unique $\lambda$, which satisfies $\|\lambda\|_{2}=1$.
\end{proposition}

Proposition~\ref{prop:dual} is proved in \S\ref{sec:setting}. It says that $\kappa(N)\|p\|_{2}$ is
\emph{how negative, on average over the odd-$\omega$ part of the band, a level-$Q$ divisor weight can be
while remaining nonnegative on the even-$\omega$ part, per unit $\ell^{2}$ norm of its coefficients}. We
call $\kappa$ the \emph{parity-separation constant of level-$Q$ divisor weights in the $\ell^{2}$
normalisation}. The constraint is one-sided in the parity: nonnegativity is imposed on the even-$\omega$
part only. The normalisation is $\ell^{2}$ in the coefficients $\lambda_{d}$, not Selberg's quadratic form
and not an $\ell^{1}$ norm; the feasible cone is the same for every norm, the constant is not. And the
level is fixed, so for each $N$ the quantity \eqref{eq:kappa} is a finite-dimensional conic distance.

\subsection{Orientation}\label{sec:orientation}
The condition $\theta_{\lambda}\ge0$ on $\vB$ is a one-sided sign condition on a divisor sum of level $Q$,
of the shape that appears in sieve theory; it is \emph{not} by itself an admissibility condition for a
classical upper-bound sieve weight, since the usual normalisation and majorisation requirements are absent.
The comparison we draw is formal: $\kappa(N)$ is the extremal value of a finite convex optimisation
problem --- a nonnegative least-squares problem in the primal form \eqref{eq:kappa}, a second-order cone
problem in the dual form of Proposition~\ref{prop:dual} --- whose constraints are divisor-sum inequalities
together with a Euclidean norm bound, and the classical parity phenomenon
\cite[ch.~16]{OperaDeCribro} concerns a different object, namely what main terms a sieve can produce
without resolving the parity of $\omega$.

\subsection{Results}\label{sec:results}

\begin{theorem}\label{thm:main}
For every fixed $A>0$,
\[
  \kappa(N)\ \ll_{A}\ (\log Q)^{-A} .
\]
In particular $\kappa(N)\to0$: by \eqref{eq:kappale1} the optimal nonnegative reweighting beats the
uniform baseline by more than any fixed power of $\log Q$.
\end{theorem}

Theorem~\ref{thm:main} is the main result. It comes from one quantitative estimate carrying a parameter
$\beta$, fixed before $N$ grows, which we state next; Theorem~\ref{thm:main} follows by letting $\beta$
grow with $A$.

\begin{theorem}\label{thm:quant}
Let $\beta\ge4$ be fixed and put $Y=L^{\beta}$. For every $\eps>0$ there is an $N_{0}(\beta,\eps)$ such
that for every integer $N\ge N_{0}(\beta,\eps)$, with $Q=\lfloor\sqrt N\rfloor$ and $L=\log Q$,
\[
  \kappa(N)\ \le\ \Bigl(\frac4{\sqrt\beta}\sqrt{\frac{15}{\pi^{2}}}+\eps\Bigr)
      \frac{L^{2-\beta/2}}{\sqrt{\log L}}
   \;+\;\bigl(4\sqrt{2e^{\,e-1}}+\eps\bigr)\,L^{2}\exp\Bigl(-\frac L{4\beta\log L}\Bigr).
\]
All constants and thresholds here may depend on $\beta$, which is fixed before $N$ grows; no statement is
claimed uniformly in $\beta$.
\end{theorem}

The first term decides the rate, the second is smaller than every power of $1/L$ but carries a large
constant. Theorem~\ref{thm:main} is the case $\beta=2A+4$. Taking instead $\beta=6$ evaluates both
leading coefficients, giving a two-term bound with numerical constants and an unevaluated threshold of
validity $N_{0}(6,\eps)$ (\S\ref{sec:effectivity}); we record it as a concrete example and not as an
optimal constant.

\begin{corollary}\label{cor:beta6}
Put
\[
  c_{1}=\frac4{\sqrt6}\sqrt{\frac{15}{\pi^{2}}}=2.0131684\ldots,\qquad
  c_{2}=4\sqrt{2e^{\,e-1}}=13.3565762\ldots .
\]
For every $\eps>0$ and all sufficiently large $N$,
\[
  \kappa(N)\ \le\ \frac{c_{1}+\eps}{L\sqrt{\log L}}
   \;+\;(c_{2}+\eps)\,L^{2}\exp\Bigl(-\frac L{24\log L}\Bigr)
  \qquad\text{and}\qquad
  \kappa(N)\ \le\ \frac{2.01317+\eps}{\log Q\,\sqrt{\log\log Q}} .
\]
\end{corollary}

The displayed decimal $2.01317$ exceeds $c_{1}$, as an upper bound requires; $c_{1}$ comes from the choice
$\beta=6$. The threshold in Theorem~\ref{thm:quant} merges those
of Proposition~\ref{prop:norm}, Theorem~\ref{thm:reduction} and Theorem~\ref{thm:PC}; see
\S\ref{sec:effectivity} for what is and is not known about it.

\paragraph{Scope: what is proved, and what is not.}\label{par:scope}
Both halves are stated here once and are not repeated elsewhere.

What is proved is a statement about one explicit finite convex problem. On the level-$Q$ incidence cone
$K_{Q}$ of \eqref{eq:kappa}, nonnegative reweighting of the even-$\omega$ columns beats the uniform
baseline \eqref{eq:kappale1} by more than any fixed power of $\log Q$ (Theorem~\ref{thm:main}); the
weights that do it are explicit \eqref{eq:X}, grouping the columns by small-prime part and reproducing
the retained classes of every $Y$-smooth row exactly; and the bound is reduced to, and then proved
from, one statement about M\"obius sums over friable integers. And at seven computable sizes
$\kappa(N)$ is located in certified rational intervals, with a certified lower bound at an eighth
(\S\ref{sec:numerics}).

What is not proved: the paper is not a breakthrough on the parity barrier of
sieve theory, and it is not a statement about primes. The parity of $\omega$ splits the band and carries
the construction and the proof; the parity obstruction of sieve theory is a different statement, about
sifted sets and primes rather than about this cone, and no implication in either direction between
Theorem~\ref{thm:main} and the obstruction cited in \S\ref{sec:orientation} is claimed or used. Every
statement above is an
upper bound on $\kappa(N)$; none of them asserts that $\kappa(N)>0$, and none says anything about
whether the rate is optimal. Every bound is proved for \emph{fixed} $\beta$, and nothing is claimed
uniformly in a growing $\beta$ (Remark~\ref{rem:fixedbeta}). The theorem is effective in principle,
but no numerical threshold and no useful finite-size estimate has been evaluated --- which is not to
say that none exists (\S\ref{sec:effectivity}). \S\ref{sec:numerics} reports what is known at
computable sizes, and none of it bears on Theorem~\ref{thm:main}.

\subsection{Method}\label{sec:method}
The proof exhibits a feasible point of $K_{Q}$ close to $t$. Fix a parameter $\beta\ge4$ and put
\begin{equation}\label{eq:Yparam}
  Y=L^{\beta},\qquad a_{Y}(n)=\prod_{p\mid n,\ p\le Y}p ,
\end{equation}
taking $\beta=6$ only where a numerical constant is wanted (Corollary~\ref{cor:beta6}). Group the band by the
\emph{small-prime part} $a_{Y}$: for squarefree $Y$-smooth $a$ let $B_{a}$ and $R_{a}$ be the number of
elements of $\vB$ and of $\vR$ with $a_{Y}=a$, and set
\begin{equation}\label{eq:X}
  x_{\ell}=\frac{R_{a_{Y}(\ell)}}{R\,B_{a_{Y}(\ell)}}\quad\text{if }a_{Y}(\ell)\le\sqrt Q
  \text{ and }B_{a_{Y}(\ell)}>0;\qquad x_{\ell}=0\quad\text{otherwise}.
\end{equation}
Then $x\ge0$, so $Ax\in K_{Q}$ and $\dist_{2}(t,K_{Q})\le\|t-Ax\|_{2}$. On every $Y$-smooth row these
weights reproduce the contribution of the retained classes exactly, and the error left on that row is
exactly the contribution of the discarded odd classes (Lemma~\ref{lem:E}); that error is smaller than
every power of $1/L$, by a Rankin argument (Lemma~\ref{lem:rank}); and the rows with a prime factor
exceeding $Y$ contribute $O((Y\log Y)^{-1/2})$, because both $t_{d}$ and $(Ax)_{d}$ are $O(1/d)$ there
(Lemma~\ref{lem:rough}, with the prime-sum estimate supplied independently by
Lemma~\ref{lem:primetail}). That is
Theorem~\ref{thm:reduction}, whose single arithmetic input is

\begin{hypothesis}[PC]\label{hyp:PC}
There is an absolute $C$ with $R_{a}\le C\,B_{a}$ for every squarefree $Y$-smooth $a\le\sqrt Q$.
\end{hypothesis}

Hypothesis~\ref{hyp:PC} says that after fixing the small-prime part, the two $\omega$-parities of the rough
cofactor occur in bounded ratio. We prove it with $C=1+O((\log L)^{-2})$ (Theorem~\ref{thm:PC}). Writing
$n=ab$ with $a=a_{Y}(n)$, the cofactor $b$ runs over squarefree integers with all prime factors in $(Y,Q]$
and $b\le N/a\in[Q^{3/2},N]$, and the parity split of such $b$ is governed by
\begin{equation}\label{eq:UM}
  U(X)=\#\{b\le X:\mu^{2}(b)=1,\ Y<\Pm(b),\ \Pp(b)\le Q\},\qquad
  M_{Y,Q}(X)=\!\!\!\sum_{\substack{b\le X\\ Y<\Pm(b),\ \Pp(b)\le Q}}\!\!\!\mu(b).
\end{equation}
Hypothesis~\ref{hyp:PC} follows from $U(X)\gg XV(Y)$ together with $|M_{Y,Q}(X)|=o(U(X))$, where
$V(y)=\prod_{p\le y}(1-1/p)$. The unsigned lower bound is a one-step Buchstab peeling against a sieve
density valid on a long range (Lemmas~\ref{lem:phi} and \ref{lem:U}).

For the signed estimate we use the following elementary fact: since $Q=\lfloor\sqrt N\rfloor$,
\[
  N<(Q+1)^{2}\le q_{1}^{2},
\]
so an integer $n\le N$ has at most one prime factor exceeding $Q$, and that prime divides $n$ to the first
power. The large-prime removal is a classical splitting step; what the inequality buys at this particular
level is that the removal is exact in one pass and that the cofactor needs \emph{no} upper cutoff, since
for $p>Q$ it is automatically smaller than $q_{1}$. Thus (Lemma~\ref{lem:1LP})
\[
  M_{Y,Q}(X)=M_{Y}(X)+\sum_{Q<p\le X}M_{Y}(X/p),\qquad
  M_{Y}(z)=\!\!\sum_{b\le z,\ \Pm(b)>Y}\!\!\mu(b),
\]
with the inner sums unrestricted above, and from there nothing is required beyond the quantitative prime
number theorem (P) and the matching M\"obius bound (L) of \S\ref{sec:inputs}
(Lemma~\ref{lem:roughmob} and Theorem~\ref{thm:signed}). No two-cutoff signed theorem has
to be quoted.

Finally, $\kappa$ carries $\|p\|_{2}$ in its denominator, so an upper bound for $\kappa$ needs a
\emph{lower} bound for $\|p\|_{2}$. We prove $\|p\|_{2}\ge|p_{1}|=(1+o(1))N/(4RL^{2})$ in
\S\ref{sec:norm}, by the same large-prime removal at $d=1$ followed by an exact integral transform; the
main term comes from $\sum_{n\ge1}\mu(n)\log n/n=-1$. That section also proves $R\gg N$, which the
construction needs as well. Both statements are instances of known asymptotics for M\"obius sums over
friable integers (\S\ref{sec:prior}); we give short self-contained proofs in the single case needed.

\section{Positioning}\label{sec:prior}

\subsection{Programming formulations of sieve weights}\label{sec:lp}
The feasible set of Proposition~\ref{prop:dual} is a parity-relaxed version of the sign constraint in the
linear- and semidefinite-programming formulations of the sieve. Brady \cite{BradyThesis,BradySDP} is the
principal comparison: the thesis formulates the choice of combinatorial sieve weights as a linear program
with pointwise divisor-sum sign constraints and studies its duality
(\cite[\S3.1, \S4.1, \S7.1]{BradyThesis}) and its optima in a model problem
(\cite[\S4.2]{BradyThesis}), computes numerical optima at finite level
(\cite[\S8.1]{BradyThesis}) and reproduces Selberg's parity example with quantitative two-term
asymptotics (\cite[\S8.2]{BradyThesis}), while
\cite{BradySDP} lifts the framework to semidefinite programming over matrix cones and shows that verifying
divisor-sum sign constraints is co-NP-hard (\cite[\S2, Prop.~1]{BradySDP}), stating that it proves no new
sieve bound. Our constraint set differs in two ways: the sign condition is imposed on the even-$\omega$
part of the band only, and the duality is over the cone generated by divisor-incidence columns rather than
over partition cones of positive semidefinite matrices. We prove an asymptotic statement about an explicitly
listed family of finite constraints and make no complexity claim; nothing here bears on \cite{BradySDP}, whose
input representation is different. Related formulations appear in Heath-Brown \cite[\S4]{HeathBrown}, where
divisor-sum weights are presented as a linear program with pointwise inequalities and a dual read off from
extremal sequences graded by the parity of $\Omega$, and in Ford--Maynard
\cite[Thms.~7.3--7.4]{FordMaynard}. Rigorous numerical work on sieve constants is established practice
(\cite[\S7.1]{Polymath}, \cite[\S\S5.1--5.2]{BookerBrowning}, and, with proved approximation-error bounds
rather than certificates, \cite[\S\S7--8]{Franze}); we claim no priority for a
programming formulation of sieve weights, for conic duality in sieve theory, or for certified sieve
computation.

\subsection{M\"obius sums over friable and over rough integers}\label{sec:friable}
The counting lemmas of \S\ref{sec:counting} and the asymptotic of \S\ref{sec:norm} are weak forms of
studied objects; we record this so that no novelty is read into them.

\begin{itemize}
\item $U(X)$ in \eqref{eq:UM} is the squarefree analogue of the two-cutoff count
$\Theta(x,y,z)=\#\{n\le x:\text{all prime factors in }(z,y]\}$. That count was first studied by
Friedlander \cite{Friedlander}, who for fixed values of the two ratios of logarithms obtained an
asymptotic whose factor is defined by a system of differential-difference equations; sharper estimates are due
to Saias \cite{SaiasI}, continued in \cite{SaiasII} (for this ordering see
\cite[\S7]{HildebrandTenenbaum}). All of them count unrestricted integers. For the unrestricted analogue, specialising the modulus
in \cite[(1.4)]{TenenbaumCribler} to the primorial through $Y$ compares the sifted friable count with
$\Psi(x,y)$, whose size is given by \cite[Thm.~1.1]{HildebrandTenenbaum}; passing to squarefree integers
requires a further adjustment. Lemma~\ref{lem:U} is a weak consequence.
\item Lemma~\ref{lem:psi} is a weak form of \cite[Cor.~1.3]{HildebrandTenenbaum}, which gives
$\Psi(x,y)=x\,u^{-(1+o(1))u}$ as $y$ and $u=\log x/\log y$ tend to infinity, uniformly for
$u\le y^{1-\eps}$ --- a range strictly containing the one used here.
\item Lemma~\ref{lem:phi} is the Buchstab asymptotic restricted to $u\to\infty$; see
\cite[ch.~III.6]{TenenbaumBook}.
\item The convolution identity behind Lemma~\ref{lem:roughmob}, and its combination with the prime
number theorem, are proved here.  Their object --- the M\"obius sum over integers free of small prime
factors --- is the subject of Alladi \cite{AlladiJNT}, which is cited as a whole for that object and
for the tradition this argument belongs to; no locator is printed and no claim is made about where
inside that paper a particular method appears. His companion paper \cite{AlladiTAMS} treats the M\"obius sum over
\emph{friable} integers, the object of \S\ref{sec:norm}; the two papers are easily conflated.
\item The asymptotic $P_{1}=-N/(\log N)^{2}+O(N/(\log N)^{3})$ of \S\ref{sec:norm} is such a friable
M\"obius sum, at $u=\log N/\log Q\to2$. For it in the form used here see de la Bret\`eche and Tenenbaum
\cite[Cor.~1.3, (1.23)]{BTWirsing}: specialising to $\mu$ and taking $y=\sqrt N+\tfrac12$ puts $u$ below
$2$, which satisfies the one-sided exclusion condition there, and gives leading term
$x\,\varpi'(u)/\log^{2}y$ with $\varpi'(u)=-1/u^{2}$, i.e.\ $-N/\log^{2}N$; changing the cutoff from $y$
to $Q$ affects at most one prime and costs $O(\sqrt N)$, as does the lower band edge. Our
\S\ref{sec:norm} is an elementary derivation in the single case needed, not a new asymptotic.
\end{itemize}

\subsection{What is new}\label{sec:new}
The construction \eqref{eq:X}: grouping the even-$\omega$ columns by small-prime part and weighting each
class by the ratio of its odd to its even count produces a feasible point of $K_{Q}$ which, on every
$Y$-smooth row, reproduces the contribution of the retained classes exactly, the remaining error being
exactly the contribution of the discarded odd classes on that row; this reduces $\kappa(N)\to0$ to the
arithmetic statement Hypothesis~\ref{hyp:PC}. Applying
this conditional class reweighting to this particular finite cone is the contribution. The large-prime
removal of Lemma~\ref{lem:1LP} is classical; the observation we use is only that at the level
$Q=\lfloor\sqrt N\rfloor$ it needs no second pass and leaves the cofactor unrestricted above, which is why
Theorem~\ref{thm:signed} needs nothing beyond the quantitative prime number theorem (P) and the matching
M\"obius bound (L). We have not located a published
form of Theorem~\ref{thm:signed} in the uniformity used here, and we claim no novelty for it: the
surrounding machinery is standard and we have not settled the question.

What the paper does and does not claim is stated once in \S\ref{par:scope} and is not restated here.

\section{Notation, inputs, and elementary lemmas}\label{sec:setting}

In addition to \eqref{eq:param}--\eqref{eq:Yparam} we write
\begin{equation}\label{eq:more}
  \theta=\frac1{\log Y},\qquad
  Z_{0}=\exp\Bigl(\frac L{(\log L)^{3}}\Bigr),\qquad
  S=\sum_{p\le Y}\frac1p,\qquad V(y)=\prod_{p\le y}\Bigl(1-\frac1p\Bigr),
\end{equation}
and, for $z\ge1$,
\begin{equation}\label{eq:counts}
  \Phi(z)=\#\{n\le z:\Pm(n)>Y\},\qquad
  M_{Y}(z)=\!\!\sum_{b\le z,\ \Pm(b)>Y}\!\!\mu(b),\qquad
  M(x)=\sum_{n\le x}\mu(n),
\end{equation}
with $n=1$ counted in $\Phi$; also $\Psi(z,y)=\#\{n\le z:\Pp(n)\le y\}$ and
$Q_{2}(x)=\sum_{n\le x}\mu^{2}(n)$, and $\gamma$ is Euler's constant. Note that $M(x)=M(\lfloor x\rfloor)$
for real $x$.

\paragraph{Asymptotic convention.}
The parameter $\beta$ of \eqref{eq:Yparam} is fixed before $N$ grows and is treated throughout as a
constant. Every constant implied by $O$, $o$ and $\ll$ below may therefore depend on $\beta$, and on
$\eps$ wherever an $\eps$ appears, and is otherwise absolute; we write a subscript only where we wish to
draw attention to the dependence, so its absence records emphasis and not uniformity. In particular
no statement below is claimed uniformly in a growing $\beta$: letting $\beta\to\infty$ with $N$ would
require the fresh uniformity analysis described in Remark~\ref{rem:fixedbeta}, at the end of
\S\ref{sec:assembly}, which we do not carry out. With $\beta$ fixed, $S\asymp\log\log L$ and $V(Y)\sim e^{-\gamma}/(\beta\log L)$ by
Mertens.

\paragraph{Standing assumption.}
From \S\ref{sec:reduction} onward, $N$ is assumed large enough that $R>0$, $p\neq0$ and
\begin{equation}\label{eq:standing}
  2\le Y<Q ,
\end{equation}
which is possible because $Y=(\log Q)^{\beta}$. Condition \eqref{eq:standing} is used in
Lemma~\ref{lem:1LP} and nowhere stated again.

\subsection{Classical inputs}\label{sec:inputs}
Each is used only in the form stated.

\begin{itemize}
\item[(M1)] $\sum_{p\le x}(\log p)/p=\log x+O(1)$; explicitly $<\log x$ for $x>1$
  \cite[(3.24)]{RosserSchoenfeld}.
\item[(M2)] $\sum_{p\le x}1/p=\log\log x+O(1)$; and for $x\ge y\ge2$,
  $\sum_{y<p\le x}1/p=\log(\log x/\log y)+O(1/\log y)$.
\item[(M3)] $V(y)=e^{-\gamma}(\log y)^{-1}\bigl(1+O(1/\log y)\bigr)$.
\item[(C)] Chebyshev: $\pi(x)\ll x/\log x$; explicitly $\pi(y)\le1.26\,y/\log y$ for $y>1$, a rounding of
  the bound $\pi(y)<1.25506\,y/\log y$ of \cite[(3.6)]{RosserSchoenfeld}.
\item[(P)] There are absolute effective $C_{\pi},c_{\pi}>0$ with
  $|\pi(y)-\li(y)|\le C_{\pi}y\exp(-c_{\pi}\sqrt{\log y})$ for $y\ge3$.
\item[(L)] There are absolute effective $C_{M},c_{M}>0$ with
  $|M(x)|\le C_{M}x\exp(-c_{M}(\log x)^{1/4})$ for real $x\ge2$.
\end{itemize}

(M1)--(M3) and (C) are Mertens' and Chebyshev's theorems. Inputs (P) and (L) are both consequences, by
standard contour integration, of the de la Vall\'ee Poussin zero-free region \cite[Thm.~3.8]{Titchmarsh}
together with the bounds for $\zeta'/\zeta$ and $1/\zeta$ on it \cite[Thm.~3.11]{Titchmarsh} and truncated
Perron inversion \cite[Lem.~3.12]{Titchmarsh}; both are quoted here in deliberately weakened shapes. For
explicit constants in (P), Theorem~2 of \cite{Trudgian} gives
\[
  |\pi(x)-\li(x)|\le0.2795\,x(\log x)^{-3/4}\exp\Bigl(-\sqrt{\tfrac{\log x}{6.455}}\Bigr)
  \qquad(x\ge229),
\]
whose logarithmic factor is at most $1$ for $x\ge e$ and whose exponential constant is
$1/\sqrt{6.455}=0.3935\ldots$; enlarging $C_{\pi}$ over the finite range $3\le x<229$ yields (P) with
$c_{\pi}=0.39$. Nothing below depends on the value of $c_{\pi}$, which enters only through
$e^{-c_{\pi}\sqrt L}$ being smaller than every power of $1/L$. We also use
\begin{equation}\label{eq:Q2}
  Q_{2}(x)=\frac6{\pi^{2}}x+O(\sqrt x),
\end{equation}
which follows from $Q_{2}(x)=\sum_{d\le\sqrt x}\mu(d)\lfloor x/d^{2}\rfloor
=x\sum_{d\le\sqrt x}\mu(d)d^{-2}+O(\sqrt x)$ and $\sum_{d>\sqrt x}d^{-2}\ll x^{-1/2}$.

\begin{lemma}\label{lem:mobconst}
$\displaystyle\sum_{n\ge1}\frac{\mu(n)}n=0$ and $\displaystyle\sum_{n\ge1}\frac{\mu(n)\log n}n=-1$.
\end{lemma}

\begin{proof}
By partial summation and (L), $\sum_{k>T}\mu(k)/k=-M(T)/T+\int_{T}^{\infty}M(v)v^{-2}dv\to0$ as
$T\to\infty$, and likewise with $\log k$ inserted; so both series converge. For $\sigma>1$ we have
$\sum_{n}\mu(n)n^{-\sigma}=\zeta(\sigma)^{-1}$ and, differentiating,
$\sum_{n}\mu(n)(\log n)n^{-\sigma}=\zeta'(\sigma)/\zeta(\sigma)^{2}$. Since a Dirichlet series converging
at $s=1$ is continuous from the right there (Abel's theorem for Dirichlet series), the two sums equal
$\lim_{\sigma\to1^{+}}\zeta(\sigma)^{-1}=0$ and
$\lim_{\sigma\to1^{+}}\zeta'(\sigma)/\zeta(\sigma)^{2}=-1$, the latter from
$\zeta(\sigma)=(\sigma-1)^{-1}+\gamma+O(\sigma-1)$ and $\zeta'(\sigma)=-(\sigma-1)^{-2}+O(1)$.
\end{proof}

\begin{lemma}\label{lem:elem}
For every integer $N\ge4$:
\begin{enumerate}
\item[(i)] $Q^{2}\le N<(Q+1)^{2}\le q_{1}^{2}$; consequently every $n\le N$ has at most one prime factor
exceeding $Q$, and that prime divides $n$ exactly once;
\item[(ii)] $Q/2-1\le\thr\le Q$;
\item[(iii)] $2L\le\log N<2L+2/Q$, so $L/\log N\le1/2$.
\end{enumerate}
\end{lemma}

\begin{proof}
(i) $Q=\lfloor\sqrt N\rfloor$ gives $Q+1>\sqrt N$, so $N<(Q+1)^{2}$; and $q_{1}\ge Q+1$ because $q_{1}$ is
an integer exceeding $Q$. Two primes exceeding $Q$ (equal or not) dividing $n\le N$ would force
$n\ge q_{1}^{2}>N$. (ii) $N<(Q+1)^{2}\le q_{1}(Q+1)$ gives $N/q_{1}<Q+1$, so $\thr\le Q$; and
$q_{1}\le2Q$ by Bertrand gives $\thr\ge N/(2Q)-1\ge Q/2-1$. (iii) $Q^{2}\le N<(Q+1)^{2}$ and
$\log(Q+1)<L+1/Q$.
\end{proof}

Part (i) is why Theorem~\ref{thm:main} concerns every large $N$ and not a subsequence: it replaces the
hypothesis $N=Q^{2}$ in the large-prime removals of Lemmas~\ref{lem:U}, \ref{lem:1LP} and
\ref{lem:identity}.

\begin{proof}[Proof of Proposition~\ref{prop:dual}]
$K_{Q}$ is finitely generated, hence a closed convex cone, so with
$K_{Q}^{\circ}=\{w:\langle w,v\rangle\le0\ \forall v\in K_{Q}\}$ Moreau's decomposition
(see e.g.\ \cite[Thm.~1.1]{Soltan}) gives
\[
  \dist_{2}(t,K_{Q})=\max\{\langle t,w\rangle:w\in K_{Q}^{\circ},\ \|w\|_{2}\le1\}.
\]
Put $\lambda=-w$; then $\langle w,A_{\ell}\rangle=\theta_{w}(\ell)\le0$ is exactly
$\theta_{\lambda}(\ell)\ge0$, and by \eqref{eq:tp}
\[
  \langle t,w\rangle=-\sum_{d}t_{d}\lambda_{d}=-\frac1R\sum_{d}c_{\vR}(d)\lambda_{d}
  =-\frac1R\sum_{n\in\vR}\theta_{\lambda}(n)=-\mathbb E_{\vR}\theta_{\lambda},
\]
which is the displayed equality. If $t\in K_{Q}$ the distance is $0$ and $w=0$ attains it. If
$t\notin K_{Q}$ the distance is positive and is attained at $w=(t-\Pi_{K_{Q}}t)/\|t-\Pi_{K_{Q}}t\|_{2}$.
For uniqueness, suppose $w_{1}\neq w_{2}$ both lie in $K_{Q}^{\circ}$ with $\|w_{i}\|_{2}\le1$ and
$\langle t,w_{i}\rangle=D>0$. Their midpoint $w$ lies in $K_{Q}^{\circ}$, has $\langle t,w\rangle=D$, and
has $\|w\|_{2}<1$ by strict convexity of the Euclidean ball; then $w/\|w\|_{2}\in K_{Q}^{\circ}$ is
feasible with value $D/\|w\|_{2}>D$, a contradiction. The same rescaling shows $\|\lambda\|_{2}=1$ at the
maximum.
\end{proof}

Finally we record the normalisation, proved in \S\ref{sec:norm}.

\begin{proposition}\label{prop:norm}
Put $P_{1}=\sum_{n\in\band}\mu(n)$, a quantity defined without dividing by $R$. There are absolute
effective constants such that, for $\log N$ large,
\[
  P_{1}=-\frac N{(\log N)^{2}}\Bigl(1+O\Bigl(\frac1{\log N}\Bigr)\Bigr).
\]
Consequently $P_{1}<0$ and $|\band|=\frac6{\pi^{2}}(1-\log2)N+O(N/L)$, so
$R\ge\frac12|\band|\ge0.09N>0$; hence $N/R=O(1)$, indeed
$N/R\to\pi^{2}/(3(1-\log2))=10.7213228\ldots$, and $p_{1}=P_{1}/R$ is well defined and satisfies, for
every $\eps\in(0,1)$ and all $L\ge L_{0}(\eps)$,
\[
  \|p\|_{2}\ \ge\ |p_{1}|\ \ge\ (1-\eps)\,\frac NR\cdot\frac1{4L^{2}}>0 .
\]
\end{proposition}

\section{The construction and the reduction}\label{sec:reduction}

Group the band by small-prime part: for squarefree $Y$-smooth $a$ put
$\vB(a)=\{\ell\in\vB:a_{Y}(\ell)=a\}$, $B_{a}=|\vB(a)|$, $\vR(a)=\{n\in\vR:a_{Y}(n)=a\}$,
$R_{a}=|\vR(a)|$, so that $\vB=\bigsqcup_{a}\vB(a)$ and $\vR=\bigsqcup_{a}\vR(a)$. With $x$ as in
\eqref{eq:X} set
\begin{equation}\label{eq:e}
  e_{d}=t_{d}-(Ax)_{d},\qquad\text{so}\qquad\dist_{2}(t,K_{Q})\le\|e\|_{2},
\end{equation}
and
\begin{equation}\label{eq:delta}
  \delta=\frac1R\#\{n\in\vR:a_{Y}(n)>\sqrt Q\}.
\end{equation}
For every \emph{retained} class $a\le\sqrt Q$, Hypothesis~\ref{hyp:PC} implies
$B_{a}=0\Rightarrow R_{a}=0$, since $R_{a}\le C\,B_{a}$ at $B_{a}=0$ forces $R_{a}=0$: that is the
whole of the hypothesis, one inequality. It matters here
because \eqref{eq:X} is defined for a retained class with $B_{a}=0<R_{a}$, which it weights by zero,
but then nothing reproduces that class's odd mass, which survives as the dead-class term of
\eqref{eq:Eid} below. Classes with $a>\sqrt Q$ are discarded outright and are not
constrained by Hypothesis~\ref{hyp:PC}. By Theorem~\ref{thm:PC} no retained class with $B_{a}=0<R_{a}$
occurs once $L$ is large.

\begin{theorem}\label{thm:reduction}
Assume Hypothesis~\ref{hyp:PC} with constant $C$. For every $\eps>0$ and all $L\ge L_{0}(\eps)$,
\[
  \kappa(N)\ \le\ (4+\eps)L^{2}\left[\sqrt{\frac{15}{\pi^{2}}}\,\max(1,C)
  \Bigl(\sum_{p>Y}p^{-2}\Bigr)^{1/2}+\Bigl(\frac{2\delta R}N\Bigr)^{1/2}\right].
\]
\end{theorem}

\subsection{The \texorpdfstring{$Y$}{Y}-smooth rows}\label{sec:smoothrows}

\begin{lemma}\label{lem:E}
Let $d\le Q$ be squarefree with $\Pp(d)\le Y$. Then, unconditionally,
\begin{equation}\label{eq:Eid}
  e_{d}=\frac1R\Bigl[\ \sum_{a>\sqrt Q,\ d\mid a}R_{a}\ +\ \sum_{a\le\sqrt Q,\ B_{a}=0,\ d\mid a}R_{a}\Bigr],
\end{equation}
and under Hypothesis~\ref{hyp:PC} the second sum vanishes, so
\begin{equation}\label{eq:Eprime}
  e_{d}=\frac1R\#\{n\in\vR:d\mid n,\ a_{Y}(n)>\sqrt Q\}\in[0,\min(\delta,t_{d})].
\end{equation}
In particular, \emph{under Hypothesis~\ref{hyp:PC}}, $e_{1}=\delta$: on the row $d=1$ the error is then
exactly the discarded odd mass. Unconditionally, \eqref{eq:Eid} at $d=1$ adds to that the mass of the
retained classes with $B_{a}=0$.
\end{lemma}

\begin{proof}
Every prime factor of $d$ is at most $Y$ and $d$ is squarefree, so for squarefree $\ell$ one has
$d\mid\ell$ if and only if $d\mid a_{Y}(\ell)$. Hence
$\#\{\ell\in\vB(a):d\mid\ell\}=B_{a}\mathbf1_{d\mid a}$ and, by \eqref{eq:X},
\[
  (Ax)_{d}=\frac1R\!\!\sum_{a\le\sqrt Q,\ B_{a}>0,\ d\mid a}\!\!R_{a},\qquad
  t_{d}=\frac1R\sum_{a:\,d\mid a}R_{a},
\]
the second because $\vR$ is partitioned by $a_{Y}$. Subtracting gives \eqref{eq:Eid}: the classes absent
from $(Ax)_{d}$ are those with $a>\sqrt Q$ and those with $a\le\sqrt Q$, $B_{a}=0$. For a retained class
$a\le\sqrt Q$ with $B_{a}=0$, Hypothesis~\ref{hyp:PC} forces $R_{a}\le C\cdot0=0$, so every term of the
second sum is zero and the sum vanishes (its index set need not be empty). For the range in
\eqref{eq:Eprime}:
$e_{d}\ge0$ termwise; $e_{d}\le t_{d}$ since $(Ax)_{d}\ge0$; and $e_{d}\le\delta$ since the set counted
lies in $\{n\in\vR:a_{Y}(n)>\sqrt Q\}$.
\end{proof}

Thus \eqref{eq:Eprime}, and with it all of Theorem~\ref{thm:reduction}, is conditional on
Hypothesis~\ref{hyp:PC}; no part of the reduction is hypothesis-free.

\begin{lemma}\label{lem:disc}
Assume Hypothesis~\ref{hyp:PC} and put $c_{0}=N/R$. Then
$\bigl\|e|_{\Pp(d)\le Y}\bigr\|_{2}\le\sqrt{2c_{0}\delta}$.
\end{lemma}

\begin{proof}
If $\delta=0$ then $e_{d}=0$ on every $Y$-smooth row by \eqref{eq:Eprime}; so assume $\delta>0$. Since
$c_{\vR}(d)\le\lfloor N/d\rfloor\le N/d$ we have $t_{d}\le c_{0}/d$, so $0\le e_{d}\le f(d)$ with
$f(w)=\min(\delta,c_{0}/w)$. As $f^{2}$ is nonincreasing on $(0,\infty)$,
$f(d)^{2}\le\int_{d-1}^{d}f(w)^{2}dw$, whence
\[
  \sum_{d\ge1}f(d)^{2}\le\int_{0}^{\infty}\min\Bigl(\delta^{2},\frac{c_{0}^{2}}{w^{2}}\Bigr)dw
  =\int_{0}^{c_{0}/\delta}\delta^{2}dw+\int_{c_{0}/\delta}^{\infty}\frac{c_{0}^{2}}{w^{2}}dw
  =2c_{0}\delta.\qedhere
\]
\end{proof}

\subsection{The rows with a large prime factor}\label{sec:roughrows}

The two ingredients are separated, because only the second uses Hypothesis~\ref{hyp:PC}: the prime-sum
estimate below is unconditional and is also the one quoted in \S\ref{sec:assembly}, so no
PC-conditional statement is used anywhere in the proof of Hypothesis~\ref{hyp:PC} itself.

\begin{lemma}\label{lem:primetail}
Unconditionally,
\[
  \sum_{\substack{d\ \text{squarefree}\\ \Pp(d)>Y}}d^{-2}
  \ \le\ \frac{\zeta(2)}{\zeta(4)}\sum_{p>Y}p^{-2},\qquad
  \sqrt{\frac{\zeta(2)}{\zeta(4)}}=\sqrt{\frac{15}{\pi^{2}}}=1.2328088\ldots,
\]
and $\sum_{p>Y}p^{-2}=(1+o(1))/(Y\log Y)$ as $Y\to\infty$.
\end{lemma}

\begin{proof}
Every squarefree $d$ with a prime factor exceeding $Y$ is $pm$ with $p>Y$ prime and $m$ squarefree, so
\[
  \sum_{\substack{d\ \text{squarefree}\\ \Pp(d)>Y}}d^{-2}
  \le\Bigl(\sum_{p>Y}p^{-2}\Bigr)\!\!\sum_{m\ \text{squarefree}}\!\!m^{-2}
  =\Bigl(\sum_{p>Y}p^{-2}\Bigr)\prod_{p}(1+p^{-2})=\frac{\zeta(2)}{\zeta(4)}\sum_{p>Y}p^{-2},
\]
because $1+p^{-2}=(1-p^{-4})/(1-p^{-2})$; and $\zeta(2)/\zeta(4)=15/\pi^{2}$. For the last claim, partial
summation with $\pi(t)=(1+o(1))t/\log t$ gives
$\sum_{p>Y}p^{-2}=-\pi(Y)/Y^{2}+2\int_{Y}^{\infty}\pi(t)t^{-3}dt=(1+o(1))(-1+2)/(Y\log Y)$.
\end{proof}

\begin{lemma}\label{lem:rough}
Assume Hypothesis~\ref{hyp:PC}. Then
\[
  \bigl\|e|_{\Pp(d)>Y}\bigr\|_{2}\le\max(1,C)\frac NR\sqrt{\frac{\zeta(2)}{\zeta(4)}}
  \Bigl(\sum_{p>Y}p^{-2}\Bigr)^{1/2}.
\]
\end{lemma}

\begin{proof}
Let $d\le Q$ be squarefree with a prime factor exceeding $Y$. Then $t_{d}=c_{\vR}(d)/R\le(N/R)/d$; and by
\eqref{eq:X} with Hypothesis~\ref{hyp:PC} one has $x_{\ell}\le C/R$ for every $\ell$, so
$0\le(Ax)_{d}\le(C/R)\#\{\ell\in\vB:d\mid\ell\}=Cc_{\vB}(d)/R\le C(N/R)/d$. Hence
$|e_{d}|\le\max(1,C)(N/R)/d$. The factor $N/R$ here is exact and is carried as such: here it enters
linearly (it also enters Lemma~\ref{lem:disc}, under a square root), and it cancels in \S\ref{sec:assembly} against the $R/N$ that
comes from $1/\|p\|_{2}$, so the density bound of Proposition~\ref{prop:norm} is not used at this step
(Remark~\ref{rem:density}). Summing the squares over those $d$ and applying
Lemma~\ref{lem:primetail} gives the claim.
\end{proof}

\subsection{The discarded mass}\label{sec:discarded}

\begin{lemma}\label{lem:rank}
With $\theta=1/\log Y$ and $\Xi=e^{\,e-1}=5.5749415\ldots$,
\[
  \sum_{n\le N}a_{Y}(n)^{\theta}\le\Xi\,N,
  \qquad\text{hence}\qquad
  \delta\le\Xi\,\frac NR\exp\Bigl(-\frac L{2\log Y}\Bigr).
\]
\end{lemma}

\begin{proof}
Since $a_{Y}(n)$ is squarefree and $Y$-smooth, for $\theta>0$
\[
  a_{Y}(n)^{\theta}=\!\!\prod_{p\mid n,\ p\le Y}\!\!\bigl(1+(p^{\theta}-1)\bigr)
  =\sum_{d\mid a_{Y}(n)}\prod_{p\mid d}(p^{\theta}-1)
  =\sum_{\substack{d\mid n,\ \mu^{2}(d)=1\\ \Pp(d)\le Y}}\ \prod_{p\mid d}(p^{\theta}-1),
\]
because $d\mid a_{Y}(n)$ is the same as $d\mid n$ with $d$ squarefree and $Y$-smooth. All coefficients are
nonnegative. Summing over $n\le N$ and using $\lfloor N/d\rfloor\le N/d$,
\[
  \sum_{n\le N}a_{Y}(n)^{\theta}\le N\prod_{p\le Y}\Bigl(1+\frac{p^{\theta}-1}p\Bigr).
\]
For $p\le Y$ we have $0\le\theta\log p\le1$, and $e^{v}-1\le(e-1)v$ on $[0,1]$ by convexity, so by (M1)
\[
  \sum_{p\le Y}\frac{p^{\theta}-1}p\le(e-1)\theta\sum_{p\le Y}\frac{\log p}p
  \le(e-1)\theta\log Y=e-1;
\]
now use $\prod(1+u_{p})\le\exp(\sum u_{p})$. For the second display, Markov's inequality on $\vR$ with
threshold $\sqrt Q$ gives
$\delta\le R^{-1}Q^{-\theta/2}\sum_{n\in\vR}a_{Y}(n)^{\theta}\le(N/R)\,\Xi\,Q^{-\theta/2}$, and
$Q^{-\theta/2}=\exp(-L/(2\log Y))$.
\end{proof}

The divisor expansion in Lemma~\ref{lem:rank} needs no restriction on the range of $a$, no sieve bound and
no splitting of the $a$-sum.

\subsection{Proof of Theorem~\ref{thm:reduction}}\label{sec:redproof}
By \eqref{eq:e} and Lemmas~\ref{lem:disc} and \ref{lem:rough},
\[
  \dist_{2}(t,K_{Q})\le\sqrt{2(N/R)\delta}+\max(1,C)\frac NR\sqrt{\frac{15}{\pi^{2}}}
  \Bigl(\sum_{p>Y}p^{-2}\Bigr)^{1/2}.
\]
By Proposition~\ref{prop:norm}, $1/\|p\|_{2}\le(4+\eps)L^{2}(R/N)$ for $L\ge L_{0}(\eps)$. Multiplying,
the factor $N/R$ cancels exactly in the second term and halves in the first. \qed

\section{The counting lemmas}\label{sec:counting}

The range that matters is $X\in[Q^{3/2},N]$, i.e.\ $\log X\in[1.5L,2L+O(1/Q)]$ by
Lemma~\ref{lem:elem}(iii). A single $Z_{0}=\exp(L/(\log L)^{3})$ serves every lemma below.

\begin{lemma}\label{lem:phi}
Uniformly for $z\ge Z_{0}$, $\ \Phi(z)=zV(Y)(1+o(1))$.
\end{lemma}

\begin{proof}
Let $P=\prod_{p\le Y}p$, $J=\lceil\log L\rceil$, and for $n\ge1$ put $k=\omega(\gcd(n,P))$. Legendre's
identity reads $\mathbf1_{\Pm(n)>Y}=\sum_{d\mid\gcd(n,P)}\mu(d)$, and
\[
  \sum_{j=0}^{J}(-1)^{j}\binom kj=(-1)^{J}\binom{k-1}J\ \ (k\ge1),\qquad=1\ \ (k=0),
\]
so the truncation error at each $n$ is $\binom{k-1}J\le\binom kJ$. Summing over $n\le z$ and using
$\sum_{n\le z}\binom{\omega(\gcd(n,P))}J=\sum_{d\mid P,\ \omega(d)=J}\lfloor z/d\rfloor$,
\begin{equation}\label{eq:B1}
  \Bigl|\Phi(z)-\!\!\!\sum_{d\mid P,\ \omega(d)\le J}\!\!\!\mu(d)\lfloor z/d\rfloor\Bigr|
  \le z\!\!\!\sum_{d\mid P,\ \omega(d)=J}\!\!\!\frac1d\ \le\ \frac{zS^{J}}{J!},
\end{equation}
the last step because the elementary symmetric function obeys $e_{J}(1/p:p\le Y)\le S^{J}/J!$. Replacing
$\lfloor z/d\rfloor$ by $z/d$ and completing the M\"obius sum,
\begin{align}
  \sum_{d\mid P,\ \omega(d)\le J}\!\!\!\mu(d)\lfloor z/d\rfloor
   &=zV(Y)-z\!\!\!\sum_{d\mid P,\ \omega(d)>J}\!\!\!\frac{\mu(d)}d
     +O\bigl(\#\{d\mid P:\omega(d)\le J\}\bigr),\label{eq:B2a}\\
  \Bigl|\sum_{d\mid P,\ \omega(d)>J}\frac{\mu(d)}d\Bigr|
   &\le\sum_{j>J}\frac{S^{j}}{j!}\le\frac{S^{J+1}}{(J+1)!}\Bigl(1-\frac S{J+2}\Bigr)^{-1}
     \le\frac{2S^{J+1}}{(J+1)!}\quad\text{if }S\le\frac{J+2}2,\label{eq:B2}\\
  \#\{d\mid P:\omega(d)\le J\}&\le(J+1)\pi(Y)^{J}\le Y^{J}.\label{eq:B3}
\end{align}
The side condition in \eqref{eq:B2} holds eventually because $S\asymp\log\log L$ and $J\asymp\log L$; the
second inequality of \eqref{eq:B3} holds eventually because (C) gives $Y/\pi(Y)\gg\log Y$, so
$(Y/\pi(Y))^{J}\ge J+1$. Divide the three errors by $zV(Y)$, which is $\asymp z/(\beta\log L)$ by (M3).
Since $S=(1+o(1))\log\log L$ and $J=\lceil\log L\rceil$,
\[
  \frac{S^{J}}{J!}\le\Bigl(\frac{eS}J\Bigr)^{J}
  =\exp\bigl(-(\log L)(\log\log L-\log\log\log L-1+o(1))\bigr),
\]
smaller than every fixed power of $1/L$; so \eqref{eq:B1} and \eqref{eq:B2} are $o(V(Y))$ with room. For
\eqref{eq:B3} we compare logarithms: $\log(Y^{J})=\beta J\log L=\beta(\log L)^{2}(1+o(1))$, whereas
\[
  \log\bigl(zV(Y)\bigr)\ \ge\ \log\bigl(Z_{0}V(Y)\bigr)
  =\frac L{(\log L)^{3}}-\log(\beta\log L)-\gamma+o(1),
\]
and $\beta(\log L)^{2}=o\bigl(L/(\log L)^{3}\bigr)$ because $(\log L)^{5}=o(L)$. All three errors are
therefore $o(1)$ relative to $zV(Y)$, uniformly in $z\ge Z_{0}$.
\end{proof}

Lemma~\ref{lem:phi} asks only for $z\ge Z_{0}$, where $u=\log z/\log Y\ge L/(\beta(\log L)^{4})\to\infty$;
this avoids the Buchstab transition rather than estimating it. For comparison, with $\varpi$ Buchstab's
function one has, uniformly for $2\le Y\le\sqrt z$,
\[
  \Phi(z)=\frac z{\log Y}\Bigl\{\varpi(u)+O\Bigl(\frac1{\log Y}\Bigr)\Bigr\}
\]
\cite[ch.~III.6]{TenenbaumBook}, hence $\Phi(z)\sim\varpi(u)e^{\gamma}zV(Y)$ as $Y\to\infty$, uniformly in
that range --- which does contain the range of Lemma~\ref{lem:phi}, since $2\beta(\log L)^{4}\le L$
eventually. Here $\max_{2\le u\le3}\varpi(u)e^{\gamma}=1.0101232\ldots$, attained at
$u=2.7632228\ldots$, and the maximum can be read off the delay equation: on $[2,3]$ one has
$u\varpi(u)=1+\log(u-1)$, so $u^{2}\varpi'(u)=u/(u-1)-1-\log(u-1)$, which written in $v=u-1$ is
$1/v-\log v$, strictly decreasing on $[1,2]$ from $1$ to $-0.1931\ldots$ and therefore zero exactly once;
there $v\log v=1$ and $\varpi(u)=1/v=0.5671432\ldots$. We make no claim about where the maximum over all
$u\ge2$ lies, because nothing below needs one. What is used is only the limit
$\varpi(u)e^{\gamma}\to1$ as $u\to\infty$, which yields the constant $1-\log2$ in Lemma~\ref{lem:U}.

\begin{lemma}\label{lem:U}
Uniformly for $X\in[Q^{3/2},N]$,
\[
  U(X)\ \ge\ (1-\log2+o(1))\,X\,V(Y),\qquad 1-\log2=0.3068528\ldots,
\]
and in particular $U(X)\gg X/\log Y$.
\end{lemma}

\begin{proof}
Let $\Phi_{Q}(z)=\#\{n\le z:Y<\Pm(n),\ \Pp(n)\le Q\}$. If $n\le X\le N$ has $\Pm(n)>Y$ and $\Pp(n)>Q$,
then by Lemma~\ref{lem:elem}(i) the large prime $p$ is unique and divides $n$ once; writing $n=pm$,
Lemma~\ref{lem:elem}(i)--(ii) give $m\le X/p\le N/q_{1}<q_{1}$, so every prime factor of $m$ is smaller
than $q_{1}$, i.e.\ at most $Q$ --- the upper cutoff on $m$ is automatic --- and $\Pm(m)>Y$ or $m=1$. The
map $n\mapsto(p,m)$ is a bijection onto the pairs counted, so
\begin{equation}\label{eq:peel}
  \Phi(X)=\Phi_{Q}(X)+\sum_{Q<p\le X}\Phi_{Q}(X/p),\quad\text{hence}\quad
  \Phi_{Q}(X)\ge\Phi(X)-\sum_{Q<p\le X}\Phi(X/p).
\end{equation}
Split the subtracted sum at $X/p=Z_{0}$, using Lemma~\ref{lem:phi} for $p\le X/Z_{0}$ and $\Phi(z)\le z$
beyond:
\[
  \sum_{Q<p\le X}\Phi(X/p)\le(1+o(1))XV(Y)\!\!\sum_{Q<p\le X}\!\!\frac1p
   +X\!\!\sum_{X/Z_{0}<p\le X}\!\!\frac1p.
\]
By (M2) and Lemma~\ref{lem:elem}(iii), $\sum_{Q<p\le X}1/p=\log(\log X/\log Q)+O(1/L)\le\log2+o(1)$,
uniformly for $X\le N$. For the short block put $\tau_{X}=\log Z_{0}/\log X=L/((\log L)^{3}\log X)$; then
\begin{equation}\label{eq:short}
  \sum_{X/Z_{0}<p\le X}\frac1p=-\log(1-\tau_{X})+O(1/L)=\tau_{X}+O(\tau_{X}^{2}+1/L)
  \le\Bigl(\frac23+o(1)\Bigr)(\log L)^{-3},
\end{equation}
because $\log X\ge1.5L$ gives $\tau_{X}\le\frac23(\log L)^{-3}$, and $1/L=o((\log L)^{-3})$. Note that
\eqref{eq:short} is an inequality, not an equality: $\tau_{X}(\log L)^{3}$ equals $2/3$ at $X=Q^{3/2}$ and
$\frac12+O(1/(QL))$ at $X=N$. Since $(\log L)^{-3}=o(V(Y))$ the short block is $o(XV(Y))$, and
$\Phi(X)=(1+o(1))XV(Y)$ by Lemma~\ref{lem:phi}; hence $\Phi_{Q}(X)\ge(1-\log2+o(1))XV(Y)$. Finally,
deleting the non-squarefree $b$ costs at most $\sum_{p>Y}X/p^{2}\ll X/(Y\log Y)=o(XV(Y))$ by
Lemma~\ref{lem:primetail}, which is unconditional.
\end{proof}

\begin{lemma}\label{lem:psi}
Let $u=\log z/\log Y$. Uniformly for $Z_{0}\le z\le N$, so that
$u\ge L/(\beta(\log L)^{4})\to\infty$,
\[
  \Psi(z,Y)\le z\exp\Bigl(-\frac12u\log u\Bigr).
\]
\end{lemma}

\begin{proof}
Rankin's trick: for $\alpha\in(0,1]$ one has
$\Psi(z,Y)\le z^{\alpha}\prod_{p\le Y}(1-p^{-\alpha})^{-1}$, since
$\sum_{\Pp(n)\le Y}(z/n)^{\alpha}\ge\Psi(z,Y)$ termwise. Take $\alpha=1-(\log u)/(\log Y)$, so that
$z^{\alpha}=z\exp(-u\log u)$ exactly. Throughout the range $\log u=(1+o(1))\log L$ and
$\log Y=\beta\log L$, so $1-\alpha=(1+o(1))/\beta$ and $\alpha\to1-1/\beta\in[3/4,1)$, uniformly. The
bound $\alpha\ge3/4$ is moreover an inequality and not only a limit, for $L$ beyond an absolute point:
$z\le N$ gives $\log z<2L+2/Q\le3L$, so $u=\log z/\log Y\le3L/(\beta\log L)\le L$, whence
$\log u\le\log L$ and
\[
  \alpha=1-\frac{\log u}{\beta\log L}\ \ge\ 1-\frac1\beta\ \ge\ \frac34
\]
since $\beta\ge4$. Estimate
$\sum_{p\le Y}p^{-\alpha}$ by Abel summation:
\[
  \sum_{p\le Y}p^{-\alpha}=Y^{-\alpha}\pi(Y)+\alpha\int_{2}^{Y}\pi(t)t^{-\alpha-1}dt .
\]
Since $Y^{1-\alpha}=u$ and $(1-\alpha)\log Y=\log u$, the boundary term is, by (P),
\[
  Y^{-\alpha}\pi(Y)=(1+o(1))\frac{Y^{1-\alpha}}{\log Y}=(1+o(1))(1-\alpha)\frac u{\log u}
  =\Bigl(\frac1\beta+o(1)\Bigr)\frac u{\log u},
\]
so it is a definite fraction of the total and not negligible. In the integral, the part $t\le\sqrt Y$ is at
most $\int_{2}^{\sqrt Y}t^{-\alpha}dt\le Y^{(1-\alpha)/2}/(1-\alpha)=O(u^{1/2}\log L)=o(u/\log u)$; on
$[\sqrt Y,Y]$, (P) gives $\pi(t)=(1+o(1))t/\log t$ uniformly, and substituting $t=Y^{s}$, so that
$t^{1-\alpha}=u^{s}$ and $dt/(t\log t)=ds/s$,
\[
  \alpha\int_{\sqrt Y}^{Y}\pi(t)t^{-\alpha-1}dt
  =(1+o(1))\,\alpha\int_{1/2}^{1}u^{s}\frac{ds}s=(1+o(1))\,\alpha\,\frac u{\log u}
  =\Bigl(1-\frac1\beta+o(1)\Bigr)\frac u{\log u}
\]
by Laplace's method, the mass sitting at $s$ near $1$, i.e.\ at $p$ near $Y$, which is where (P) is
strongest. The two contributions sum to $(1+o(1))u/\log u$. For the Euler product, peel off $p=2$, whose
factor contributes $-\log(1-2^{-\alpha})=O(1)$ because $2^{-\alpha}\le2^{-3/4}=0.5946\ldots<1$; for
$p\ge3$ and $\alpha\ge3/4$ one has $p^{-\alpha}\le3^{-3/4}=0.4386\ldots<1/2$, so $-\log(1-v)\le v+v^{2}$
applies on its stated range, and $\sum_{3\le p\le Y}p^{-2\alpha}\le\sum_{p}p^{-3/2}=O(1)$. Hence
$\log\prod_{p\le Y}(1-p^{-\alpha})^{-1}=(1+o(1))u/\log u+O(1)$, and
$\log(\Psi(z,Y)/z)\le-u\log u+(1+o(1))u/\log u+O(1)\le-\frac12u\log u$, since $u/\log u=o(u\log u)$.
\end{proof}

\section{The parity hypothesis is a theorem}\label{sec:PC}

\subsection{The parity split}\label{sec:split}
Fix a squarefree $Y$-smooth $a\le\sqrt Q$. Every $n\in\band$ with $a_{Y}(n)=a$ is $n=ab$ with $b$
squarefree, $Y<\Pm(b)$, $\Pp(b)\le Q$ and $b\in(\thr/a,N/a]$; coprimality of $a$ and $b$ is automatic from
$\Pm(b)>Y\ge\Pp(a)$, and $\omega(n)=\omega(a)+\omega(b)$. With
\[
  U_{a}=U(N/a)-U(\thr/a),\qquad M_{a}=M_{Y,Q}(N/a)-M_{Y,Q}(\thr/a)
\]
in the notation \eqref{eq:UM}, and $\mu(b)=(-1)^{\omega(b)}$ for squarefree $b$,
\begin{equation}\label{eq:par}
  \{B_{a},R_{a}\}=\Bigl\{\frac{U_{a}+M_{a}}2,\ \frac{U_{a}-M_{a}}2\Bigr\},\qquad
  \frac{R_{a}}{B_{a}}\le\frac{U_{a}+|M_{a}|}{U_{a}-|M_{a}|}\quad\text{if }|M_{a}|<U_{a},
\end{equation}
which of the two is which being decided by the parity of $\omega(a)$. So Hypothesis~\ref{hyp:PC} needs only
$|M_{a}|$ bounded away from $U_{a}$ by a fixed factor; we prove the stronger $|M_{a}|=o(U_{a})$ because it
is available.

\subsection{The signed estimate}\label{sec:signed}

\begin{lemma}\label{lem:1LP}
Assume \eqref{eq:standing}. Then for every $X\le N$,
\[
  M_{Y,Q}(X)=M_{Y}(X)+\sum_{Q<p\le X}M_{Y}(X/p).
\]
\end{lemma}

\begin{proof}
Split $M_{Y}(X)$ according to whether $\Pp(b)\le Q$. If not, $\Pp(b)=p>Q$; by Lemma~\ref{lem:elem}(i) $p$
is the only prime factor of $b$ exceeding $Q$ and divides $b$ once, so $b=pb'$ with
$b'=b/p\le X/p\le N/q_{1}<q_{1}$, whence every prime factor of $b'$ is at most $Q$, $p\nmid b'$, and
$\Pm(b')>Y$ or $b'=1$. Conversely, let $Q<p\le X$ and $b'\le X/p$ with $\Pm(b')>Y$, and put $b=pb'$. Then
$\Pp(b')\le p$ automatically, and $\Pm(b)=\min(p,\Pm(b'))>Y$ because $p>Q\ge Y$ by \eqref{eq:standing}; so
$b$ is counted by $M_{Y}(X)$ with $\Pp(b)>Q$, and the correspondence is a bijection. Since
$\mu(pb')=-\mu(b')$, we get $M_{Y}(X)=M_{Y,Q}(X)-\sum_{Q<p\le X}M_{Y}(X/p)$.
\end{proof}

The hypothesis $Y\le Q$ is needed: without it the converse direction fails. At $N=256$, $Q=16$,
$Y=(\log16)^{6}>256$ and $X=17$, both sides of the identity would read $1$ and $2$ respectively, the
discrepancy being the pair $(p,b')=(17,1)$, whose product $17$ is not counted by $M_{Y}(17)$.

\begin{lemma}\label{lem:roughmob}
Uniformly for $Z_{0}\le z\le N$ and for every fixed $A_{0}>0$,
$\ |M_{Y}(z)|\ll_{A_{0},\beta}z(\log L)^{-A_{0}}$.
\end{lemma}

\begin{proof}
Every $n$ factors uniquely as $n=mb$ with $m$ $Y$-friable and $\Pm(b)>Y$, and $\mu(n)=\mu(m)\mu(b)$; in
Dirichlet series,
$\sum_{\Pm(b)>Y}\mu(b)b^{-s}=\prod_{p>Y}(1-p^{-s})=\zeta(s)^{-1}\prod_{p\le Y}(1-p^{-s})^{-1}$, so
$M_{Y}(z)=\sum_{m\le z,\ \Pp(m)\le Y}M(z/m)$. Split at $m\le z/W$ with $W=\exp(\sqrt{\log z})$. For
$m\le z/W$, apply (L) at argument at least $W$ and use
$\sum_{m\ Y\text{-friable}}1/m=V(Y)^{-1}\asymp e^{\gamma}\log Y$ from (M3):
\[
  \sum_{m\le z/W}|M(z/m)|\ll z\exp\bigl(-c_{M}(\log W)^{1/4}\bigr)V(Y)^{-1}
  \ll z(\log L)\exp\bigl(-c_{M}(\log z)^{1/8}\bigr).
\]
For $m>z/W$ use $|M(z/m)|\le z/m\le W$ and Lemma~\ref{lem:psi}: that block is at most
$W\Psi(z,Y)\le z\exp(\sqrt{\log z}-\frac12u\log u)$. In the stated range $\log u=(1+o(1))\log L$, so
\[
  u\log u=\frac{\log z}{\beta\log L}(1+o(1))\log L=(1+o(1))\frac{\log z}\beta,
\]
and $\sqrt{\log z}=o(\log z)$; hence $\sqrt{\log z}-\frac12u\log u\le-(\log z)/(3\beta)$ for $L$ large,
uniformly. Since $\log z\ge L/(\log L)^{3}$, both blocks are $\ll z\exp(-c'L^{1/9})$ with $c'$ depending on
$\beta$, which beats $z(\log L)^{-A_{0}}$ for every fixed $A_{0}$.
\end{proof}

\begin{theorem}\label{thm:signed}
Uniformly for $X\in[Q^{3/2},N]$,
\[
  |M_{Y,Q}(X)|\le\Bigl(\frac23+o(1)\Bigr)X(\log L)^{-3}.
\]
\end{theorem}

\begin{proof}
Apply Lemma~\ref{lem:1LP} and split the prime sum at $X/p=Z_{0}$. First
$|M_{Y}(X)|\ll X(\log L)^{-4}$ by Lemma~\ref{lem:roughmob}, as $X\ge Q^{3/2}\gg Z_{0}$. Next, for
$Q<p\le X/Z_{0}$, Lemma~\ref{lem:roughmob} with $A_{0}=4$ gives $|M_{Y}(X/p)|\ll(X/p)(\log L)^{-4}$, while
by (M2) with Lemma~\ref{lem:elem}(iii) one has $\sum_{Q<p\le X}1/p=\log(\log X/\log Q)+O(1/L)$
uniformly, which since $X\le N$ is at most $\log2+o(1)$, the limit at $X=Q^{3/2}$ being $\log(3/2)$;
only the resulting $O(1)$ is used, so this block
is $\ll X(\log L)^{-4}$. Finally, for $X/Z_{0}<p\le X$ use $|M_{Y}(z)|\le z$ and \eqref{eq:short}, which
gives at most $(\frac23+o(1))X(\log L)^{-3}$.
\end{proof}

\subsection{The hypothesis}\label{sec:PCproof}

\begin{theorem}\label{thm:PC}
For all $L$ beyond a point depending only on $\beta$, Hypothesis~\ref{hyp:PC} holds with
\[
  \frac{R_{a}}{B_{a}}\le1+\Bigl(\frac{4\beta e^{\gamma}}{3(1-\log2)}+o(1)\Bigr)(\log L)^{-2}
  =1+O_{\beta}\bigl((\log L)^{-2}\bigr),
\]
uniformly in squarefree $Y$-smooth $a\le\sqrt Q$. At $\beta=6$ the displayed constant is
$46.4345\ldots$. In particular $B_{a}>0$ for every such $a$. The implied constant is proportional to
$\beta$, as the displayed coefficient shows; the uniformity asserted is in $a$, at a fixed $\beta$.
\end{theorem}

\begin{proof}
For $a\le\sqrt Q$ we have $X:=N/a\in[N/\sqrt Q,N]\subseteq[Q^{3/2},N]$, because $N\ge Q^{2}$ gives
$N/\sqrt Q\ge Q^{3/2}$. The lower band edge costs nothing, uniformly: by Lemma~\ref{lem:elem}(ii),
$\thr/a\le Q/a$, so $U(\thr/a)\le Q/a$, $|M_{Y,Q}(\thr/a)|\le Q/a$, and
$(Q/a)/\bigl((N/a)V(Y)\bigr)=Q/(NV(Y))\ll(\log L)/Q\to0$. Hence by Lemma~\ref{lem:U},
Theorem~\ref{thm:signed} and $V(Y)=(1+o(1))e^{-\gamma}/(\beta\log L)$ from (M3),
\[
  U_{a}\ge(1-\log2+o(1))\frac NaV(Y),\qquad
  |M_{a}|\le\Bigl(\frac23+o(1)\Bigr)\frac Na(\log L)^{-3},
\]
\[
  r:=\frac{|M_{a}|}{U_{a}}\le\frac{(\frac23+o(1))(\log L)^{-3}}{(1-\log2+o(1))V(Y)}
  =\frac{\frac23\beta e^{\gamma}+o(1)}{(1-\log2)(\log L)^{2}} .
\]
So $|M_{a}|<U_{a}$ for $L$ large, whence $B_{a}\ge(U_{a}-|M_{a}|)/2>0$, and \eqref{eq:par} gives
$R_{a}/B_{a}\le(1+r)/(1-r)=1+2r+O(r^{2})$, which is the display. Every estimate depends on $a$ only
through $X=N/a$, and Lemma~\ref{lem:U} and Theorem~\ref{thm:signed} are uniform on $[Q^{3/2},N]$, so no
union bound over the classes is required.
\end{proof}

\section{The normalisation}\label{sec:norm}

This section proves Proposition~\ref{prop:norm}. Write $Y_{1}=\thr$, $P_{1}=\sum_{n\in\band}\mu(n)$,
\[
  \Phi_{\mu}(x)=\!\!\sum_{m\le x,\ \Pp(m)\le Q}\!\!\mu(m),\qquad
  F(y)=\#\{n\le y:\mu^{2}(n)=1,\ \Pp(n)\le Q\},\qquad \Delta_{\pi}=\pi-\li .
\]
Nothing in this section involves $Y$ or the construction.

\begin{lemma}\label{lem:identity}
For every integer $N\ge4$,
\[
  P_{1}=M(N)-M(Y_{1})\bigl(1+\pi(Q)\bigr)+\sum_{k\le Y_{1}}\mu(k)\pi(N/k).
\]
\end{lemma}

\begin{proof}
$P_{1}=\Phi_{\mu}(N)-\Phi_{\mu}(Y_{1})$, and $\Phi_{\mu}(Y_{1})=M(Y_{1})$ because $Y_{1}=\thr\le Q$ by
Lemma~\ref{lem:elem}(ii) makes every $m\le Y_{1}$ automatically $Q$-friable. For $\Phi_{\mu}(N)$: by
Lemma~\ref{lem:elem}(i) a squarefree $m\le N$ has at most one prime factor exceeding $Q$, and if $m=pm'$
with $p>Q$ then $m'\le\lfloor N/q_{1}\rfloor=Y_{1}\le Q$, so $m'$ is $Q$-friable, $p\nmid m'$ and
$\mu(m)=-\mu(m')$. Summing,
$M(N)=\Phi_{\mu}(N)-\sum_{Q<p\le N}\Phi_{\mu}(N/p)=\Phi_{\mu}(N)-\sum_{Q<p\le N}M(N/p)$. Finally
$\sum_{Q<p\le N}M(N/p)=\sum_{k}\mu(k)\#\{p:Q<p\le N/k\}=\sum_{k\le Y_{1}}\mu(k)[\pi(N/k)-\pi(Q)]$, the
terms with $k>Y_{1}$ vanishing because then $N/k<q_{1}$ and no prime lies in $(Q,q_{1})$; and
$\sum_{k\le Y_{1}}\mu(k)=M(Y_{1})$.
\end{proof}

\begin{lemma}\label{lem:transform}
With $T=\sum_{k\le Y_{1}}\mu(k)\li(N/k)$,
\[
  T=M(Y_{1})\li(N/Y_{1})+N\int_{1}^{Y_{1}}\frac{M(v)}{v^{2}\log(N/v)}\,dv .
\]
\end{lemma}

\begin{proof}
Write $\li(N/k)=\li(2)+\Li(N/k)$ with $\Li(y)=\int_{2}^{y}dt/\log t$, legitimate since
$N/k\ge N/Y_{1}\ge q_{1}>Q\ge2$. Then
\[
  \sum_{k\le Y_{1}}\mu(k)\int_{2}^{N/k}\frac{dt}{\log t}
  =\int_{2}^{N}\frac{M\bigl(\min(Y_{1},N/t)\bigr)}{\log t}dt
  =M(Y_{1})\Li(N/Y_{1})+\int_{N/Y_{1}}^{N}\frac{M(N/t)}{\log t}dt,
\]
the split being legitimate because $N/Y_{1}\ge q_{1}>2$; and $v=N/t$ turns the last integral into
$N\int_{1}^{Y_{1}}M(v)v^{-2}\log(N/v)^{-1}dv$. Adding $\li(2)M(Y_{1})$ and using $\li(2)+\Li(y)=\li(y)$
gives the statement.
\end{proof}

\begin{lemma}\label{lem:moments}
Put $J_{i}=\int_{1}^{Y_{1}}M(v)(\log v)^{i}v^{-2}dv$ and
$m_{i}=\sum_{k\le Y_{1}}\mu(k)(\log k)^{i}/k$. Then
\begin{enumerate}
\item[(i)] $|J_{i}|\le\Lambda_{i}=O_{i}(1)$ for every $i\ge0$, uniformly in $Y_{1}$;
\item[(ii)] $J_{0}=m_{0}-M(Y_{1})/Y_{1}$ and $J_{1}=m_{1}+J_{0}-M(Y_{1})\log Y_{1}/Y_{1}$;
\item[(iii)] $m_{0}=O(\eta_{0})$ and $m_{1}=-1+O(\eta_{1})$, where
$\eta_{j}\ll(\log Y_{1})^{(4j+3)/4}\exp\bigl(-c_{M}(\log Y_{1})^{1/4}\bigr)$.
\end{enumerate}
\end{lemma}

\begin{proof}
(i) On $[1,2)$ use $|M(v)|=|M(\lfloor v\rfloor)|\le v$, so that
\[
  \int_{1}^{2}|M(v)|(\log v)^{i}\frac{dv}{v^{2}}
  \le\int_{1}^{2}(\log v)^{i}\frac{dv}v=\frac{(\log2)^{i+1}}{i+1}.
\]
On $[2,\infty)$ apply (L) at the real argument $v$, which is legitimate since $M(v)=M(\lfloor v\rfloor)$ and
(L) is stated for real $x\ge2$; then substitute $w=\log v$ and extend the lower limit from $\log2$ down
to $0$, so as to obtain
\begin{align*}
  \int_{2}^{\infty}|M(v)|(\log v)^{i}\frac{dv}{v^{2}}
   &\le C_{M}\int_{\log2}^{\infty}e^{-c_{M}w^{1/4}}w^{i}\,dw
    \le C_{M}\int_{0}^{\infty}e^{-c_{M}w^{1/4}}w^{i}\,dw\\
   &=4C_{M}(4i+3)!\,c_{M}^{-(4i+4)} .
\end{align*}
(ii) is partial summation over the unit steps of $M$, with $Y_{1}$ an integer:
$J_{0}=\sum_{k<Y_{1}}M(k)(1/k-1/(k+1))$ telescopes to $m_{0}-M(Y_{1})/Y_{1}$, and with
$a_{k}=(\log k+1)/k$, $J_{1}=\sum_{k<Y_{1}}M(k)(a_{k}-a_{k+1})=m_{1}+m_{0}-M(Y_{1})(1+\log Y_{1})/Y_{1}$,
which is the stated form. (iii) By Lemma~\ref{lem:mobconst} the completed series are $0$ and $-1$, so
$m_{j}$ differs from its limit by $-\sum_{k>Y_{1}}\mu(k)(\log k)^{j}/k$, which partial summation and (L)
bound by $\eta_{j}$ as displayed.
\end{proof}

Since $\log Y_{1}\ge\log(Q/2-1)\ge L-\log4$ for $Q\ge8$ by Lemma~\ref{lem:elem}(ii), both $\eta_{j}$ are
$\ll L^{(4j+3)/4}\exp(-c'L^{1/4})$, smaller than every power of $1/L$.

\begin{proof}[Proof of Proposition~\ref{prop:norm}]
By Lemma~\ref{lem:identity} and $\pi=\li+\Delta_{\pi}$,
\[
  P_{1}=\underbrace{M(N)}_{\mathcal R_{1}}
   -\underbrace{M(Y_{1})(1+\pi(Q))}_{\mathcal R_{2}}
   +\underbrace{\sum_{k\le Y_{1}}\mu(k)\Delta_{\pi}(N/k)}_{\mathcal R_{3}}+T .
\]
By (L), $|\mathcal R_{1}|\le C_{M}N\exp(-c_{M}(\log N)^{1/4})$. By (L) and (C),
\[
  |\mathcal R_{2}|\le C_{M}Q\exp\bigl(-c_{M}(\log Y_{1})^{1/4}\bigr)
  \bigl(1+1.26\,Q/L\bigr)\le3C_{M}N\exp\bigl(-c'L^{1/4}\bigr)/L,
\]
using $1\le Q/L$ and $Q^{2}\le N$. For $\mathcal R_{3}$, every argument satisfies
$N/k\ge N/Y_{1}\ge q_{1}>Q$, so $\log(N/k)>L$ and (P) gives
\[
  |\mathcal R_{3}|\le C_{\pi}Ne^{-c_{\pi}\sqrt L}\!\!\sum_{k\le Y_{1}}\frac1k
  \le C_{\pi}N(1+L)e^{-c_{\pi}\sqrt L}.
\]
All three are $\ll N\exp(-c''L^{1/4})$, hence $o\bigl(N/(\log N)^{A}\bigr)$ for every $A$.

For $T$, Lemma~\ref{lem:transform} together with the elementary estimate
$\li(y)=(y/\log y)(1+O(1/\log y))$ gives
\[
  \bigl|M(Y_{1})\li(N/Y_{1})\bigr|\ll\frac NL\exp\bigl(-c'(\log Y_{1})^{1/4}\bigr),
\]
again negligible; we do not use (C) here,
which bounds $\pi$ and not $\li$. For the integral, put $w=\log v/\log N$, which lies in $[0,1/2]$ by
Lemma~\ref{lem:elem}(ii)--(iii), and expand $(1-w)^{-1}=1+w+w^{2}(1-w)^{-1}$, the last term here being at
most $2w^{2}$:
\[
  \int_{1}^{Y_{1}}\frac{M(v)\,dv}{v^{2}\log(N/v)}
  =\frac1{\log N}\Bigl[J_{0}+\frac{J_{1}}{\log N}+E\Bigr],\qquad
  |E|\le\frac{2\Lambda_{2}}{(\log N)^{2}} .
\]
By Lemma~\ref{lem:moments}(ii)--(iii), $J_{0}=O(\eta_{0}+C_{M}e^{-c'L^{1/4}})$ and
$J_{1}=-1+O(\eta_{1}+\eta_{0}+LC_{M}e^{-c'L^{1/4}})$, so
\[
  T=-\frac N{(\log N)^{2}}+O\Bigl(\frac{N\eta_{0}}{\log N}\Bigr)
   +O\Bigl(\frac{N\eta_{1}}{(\log N)^{2}}\Bigr)
   +O\Bigl(\frac{2N\Lambda_{2}}{(\log N)^{3}}\Bigr)
   +O\Bigl(\frac NLe^{-c'L^{1/4}}\Bigr).
\]
Each error is $O\bigl(N/(\log N)^{3}\bigr)$ once $\log N$ exceeds a threshold determined by
$(C_{M},c_{M},C_{\pi},c_{\pi})$ alone, which gives the displayed asymptotic. In particular $P_{1}<0$, and
for every $\eps>0$ one has $|P_{1}|\ge(1-\eps)N/(\log N)^{2}$ for $\log N$ large; with
$\log N<2L+2/Q$ from Lemma~\ref{lem:elem}(iii) this yields
$|p_{1}|=|P_{1}|/R\ge(1-\eps')(N/R)/(4L^{2})$ --- legitimate because $R>0$, which the band count below
supplies --- and $\|p\|_{2}\ge|p_{1}|$ because $d=1$ is a coordinate.

For the band size, the bijection of Lemma~\ref{lem:identity} applied to the squarefree indicator gives,
exactly, $F(N)=Q_{2}(N)-\sum_{Q<p\le N}Q_{2}(N/p)$. By \eqref{eq:Q2},
$Q_{2}(x)=\frac6{\pi^{2}}x+O(\sqrt x)$; by (P), $\sum_{Q<p\le N}1/p=\log(\log N/\log Q)+O(L^{-2})
=\log2+O(L^{-2})$; and by (C) with partial summation, $\sum_{Q<p\le N}\sqrt{N/p}\ll N/L$. Hence
$F(N)=\frac6{\pi^{2}}(1-\log2)N+O(N/L)$, and since $F(\thr)\le\thr\le Q$ we get
$|\band|=\frac6{\pi^{2}}(1-\log2)N+O(N/L)$. Finally $|\band|=|\vB|+R$ and $P_{1}=|\vB|-R$, so
$R=\frac12(|\band|-P_{1})\ge\frac12|\band|$ once $P_{1}\le0$; as
$\frac3{\pi^{2}}(1-\log2)=0.0932720\ldots$, this gives $R\ge0.09N$ for $L$ large, and
$N/R\to\pi^{2}/(3(1-\log2))$ follows from $P_{1}=o(N)$.
\end{proof}

\begin{remark}\label{rem:density}
The lower bound is at bottom the nonvanishing of one classical constant: the main term of $T$ comes from
$\sum_{n\ge1}\mu(n)\log n/n=-1$, while $\sum_{n\ge1}\mu(n)/n=0$ annihilates the would-be larger term. Only
the sign of $P_{1}$ is used for $R\gg N$. As recorded in \S\ref{sec:friable}, the asymptotic itself is not
new; we give the derivation because it is short in the single case needed and keeps the normalisation on the
same elementary footing as the rest of the proof.

It is worth saying how much of this the final bound actually consumes. The quantitative density
$N/R=O(1)$ is what makes the \emph{absolute} row-error statements of \S\ref{sec:reduction} meaningful,
$t_{d}\le(N/R)/d$ and $e_{d}\le\min(\delta,(N/R)/d)$ being bounds on quantities carrying the factor
$1/R$. In the \emph{relative} estimate it largely cancels: once Rankin is inserted, each of the two
terms of \S\ref{sec:assembly} acquires a factor $N/R$ from the numerator and a factor $R/N$ from
$1/\|p\|_{2}$, so what survives of Proposition~\ref{prop:norm} at the end is the lower bound on
$\|p\|_{2}$ and not the density. Nor is the density needed for that lower bound: the inequality
$|p_{1}|\ge(1-\eps)N/(4RL^{2})$ is the asymptotic for $P_{1}$ divided by $R$, and needs only $R>0$. So
$N/R=O(1)$ is information about the band rather than a step the final estimate cannot do without; we
record it because the absolute row bounds are stated in those terms and because the limiting value is
worth knowing, not because the argument turns on it.
\end{remark}

\section{Proof of Theorems~\ref{thm:main} and~\ref{thm:quant}}\label{sec:assembly}

Let $\beta\ge4$ be fixed. By Theorem~\ref{thm:PC}, Hypothesis~\ref{hyp:PC} holds with
$\max(1,C)=1+o(1)$. Insert into Theorem~\ref{thm:reduction} the estimate
$\sum_{p>Y}p^{-2}=(1+o(1))/(Y\log Y)=(1+o(1))/(\beta L^{\beta}\log L)$ of Lemma~\ref{lem:primetail}, and the
bound $\delta R/N\le\Xi\exp(-L/(2\beta\log L))$ of Lemma~\ref{lem:rank}, noting that $N/R$ cancels in the
latter as well. The first term of Theorem~\ref{thm:reduction} becomes
\begin{equation}\label{eq:t1}
  (4+\eps)L^{2}\sqrt{\frac{15}{\pi^{2}}}\cdot\frac{1+o(1)}{\sqrt{\beta L^{\beta}\log L}}
  =\bigl(1+o(1)\bigr)\frac{4+\eps}{\sqrt\beta}\sqrt{\frac{15}{\pi^{2}}}\cdot
   \frac{L^{2-\beta/2}}{\sqrt{\log L}},
\end{equation}
and the second becomes
\begin{equation}\label{eq:t2}
  (4+\eps)L^{2}\Bigl(2\Xi\exp\bigl(-\tfrac L{2\beta\log L}\bigr)\Bigr)^{1/2}
  =(4+\eps)\sqrt{2e^{\,e-1}}\,L^{2}\exp\Bigl(-\frac L{4\beta\log L}\Bigr),
\end{equation}
an equality rather than an asymptotic, because $\Xi=e^{\,e-1}$ exactly.

Together, \eqref{eq:t1} and \eqref{eq:t2} are Theorem~\ref{thm:quant}: absorbing the $(1+o(1))$ of
\eqref{eq:t1} into an arbitrary $\eps>0$ gives the two displayed coefficients
$4\beta^{-1/2}\sqrt{15/\pi^{2}}+\eps$ and $4\sqrt{2e^{\,e-1}}+\eps$, and every $o(1)$ and every threshold
used along the way depends only on $\beta$ and $\eps$, both fixed before $N$ grows. \qed

For Theorem~\ref{thm:main}, fix $A>0$ and apply Theorem~\ref{thm:quant} with $\beta=2A+4$. The first term
is $\ll_{A}L^{2-\beta/2}(\log L)^{-1/2}=L^{-A}(\log L)^{-1/2}\le L^{-A}$, and the second is $\ll_{A}L^{-A}$
because $\exp(-L/((8A+16)\log L))$ decays faster than every power of $L$. Hence
$\kappa(N)\ll_{A}L^{-A}=(\log Q)^{-A}$. \qed

For Corollary~\ref{cor:beta6}, take $\beta=6$: then \eqref{eq:t1} is
$(1+o(1))(4+\eps)\sqrt{15/\pi^{2}}/\sqrt6\cdot(L\sqrt{\log L})^{-1}$ and \eqref{eq:t2} is
$(4+\eps)\sqrt{2e^{\,e-1}}L^{2}\exp(-L/(24\log L))$. Since $\eps>0$ was arbitrary and
$4\sqrt{15/\pi^{2}}/\sqrt6=c_{1}$, $4\sqrt{2e^{\,e-1}}=c_{2}$, the first stated bound follows. For the
second, what matters is not that the exponential term tends to $0$ but that it is small \emph{relative
to the first term}: $L^{2}\exp(-L/(24\log L))$ decays faster than every power of $L$, hence is
$o\bigl(1/(L\sqrt{\log L})\bigr)$, so it can be absorbed into the tolerance of the first term, and since
$\eps>0$ was arbitrary the second stated bound follows as well. \qed

\begin{remark}\label{rem:fixedbeta}
The window $\beta\ge4$ is what \eqref{eq:t1} permits: at $\beta=4$ the exponent $2-\beta/2$ vanishes and
the first term is $4\sqrt{15/\pi^{2}}/(2\sqrt{\log L})=2.4656\ldots/\sqrt{\log L}$, still $o(1)$. The
constant $c_{1}$ of Corollary~\ref{cor:beta6} is thus an artefact of one choice of $\beta$, which is why
Theorem~\ref{thm:main} and not the numerical form is stated as the main result. Note also that letting $\beta$ grow with $N$
(\S\ref{par:scope}) would require a fresh uniformity analysis of Lemmas~\ref{lem:phi},
\ref{lem:psi} and \ref{lem:roughmob} and of Theorem~\ref{thm:PC}, which we do not carry out.
\end{remark}

\section{Effectivity}\label{sec:effectivity}

Every input is effective: (M1)--(M3), (C) and \eqref{eq:Q2} with explicit constants, and (P), (L) with
effective $(C_{\pi},c_{\pi})$, $(C_{M},c_{M})$ from the de la Vall\'ee Poussin zero-free region. No
Siegel-type or zero-density input is used, so Theorem~\ref{thm:main} is effective in principle. What is
not evaluated here is the threshold of validity: Corollary~\ref{cor:beta6}'s leading coefficients
$c_{1}$, $c_{2}$ are evaluated, while the constant implied by $\ll_{A}$ in Theorem~\ref{thm:main} and the
threshold beyond which any of these bounds holds are not. Two thresholds should
be kept apart. The \emph{normalisation} threshold, the point beyond which Proposition~\ref{prop:norm}
holds, is determined by $(C_{M},c_{M},C_{\pi},c_{\pi})$ through Lemma~\ref{lem:moments}, where the
explicit shape $\Lambda_{2}=O(1)+4C_{M}\cdot11!\,c_{M}^{-12}$ already appears, and we make no attempt to
optimise it. The \emph{overall} threshold $N_{0}(\beta,\eps)$ of Theorem~\ref{thm:quant} is larger: it
merges that one with the thresholds of Theorem~\ref{thm:reduction} and Theorem~\ref{thm:PC}, and it
depends in addition on the fixed $\beta$ and on the requested $\eps$, neither of which enters the
normalisation.

Two features of the argument make such an evaluation unrewarding at present, and we record them because
they bound what a computation could show.

\begin{enumerate}
\item \emph{The two-term expression of Corollary~\ref{cor:beta6} is itself larger than $1$ over a very
long range.} Evaluating that expression --- the case $\beta=6$, which is the only one with numerical
coefficients; Theorem~\ref{thm:main} has an unevaluated Vinogradov constant and nothing to substitute
into --- as if every $o(1)$ and $\eps$ were $0$, the sum
$c_{1}/(L\sqrt{\log L})+c_{2}L^{2}\exp(-L/(24\log L))$ first falls below $1$ at $L=3766.52\ldots$, i.e.\
$\log_{2}N=10867.87\ldots$. This is arithmetic on a formal expression and not a threshold for the
theorem. At that $L$ one has $Z_{0}=\exp(L/(\log L)^{3})\approx852$ against $Y=L^{6}\approx2.9\cdot10^{21}$, so
$u=\log Z_{0}/\log Y\approx0.14<1$, whereas Lemmas~\ref{lem:phi}, \ref{lem:psi} and \ref{lem:roughmob} all
require $u\to\infty$. The quantity $u$ at $z=Z_{0}$, namely $L/(6(\log L)^{4})$, first reaches $1$ near
$L=1.08\cdot10^{5}$, i.e.\ $\log_{2}N\approx3.1\cdot10^{5}$ --- a factor about $29$ in $L$ beyond the point
where the formal expression dips below $1$.
\item \emph{At every size at which we have evaluated it, the construction is empty.} The truncation in
\eqref{eq:X} keeps a class only when $a_{Y}(\ell)\le\sqrt Q$, and for $\beta=6$ the condition $Y<Q$ of
\eqref{eq:standing} is $6\log L<L$, i.e.\ $L>16.9988\ldots$ and $\log_{2}N\ge49.05$, while
$Y\le\sqrt Q$ needs $\log_{2}N\ge132.51$, both figures being the continuous crossing rounded
\emph{upward}, as a sufficient threshold requires. At $N=2^{28}$, the largest size at which
\S\ref{sec:numerics} certifies both endpoints, $Y=L^{6}\approx8.4\cdot10^{5}>Q$: every band element is $Y$-smooth, so every
small-prime part exceeds $\sqrt Q$ and $x$ vanishes identically. This is a statement about the sizes we
evaluated, not about all computable $N$. Indeed the construction is already nonempty a little past
$\log_{2}N=49.05$, and a witness can be exhibited without building a band: at $N=2^{52}$ one has
$Q=2^{26}=67108864$, $L=18.0218\ldots$, $Y=L^{6}=34260433.5\ldots<Q$, $q_{1}=67108879$ and
$\thr=67108849$, and the primes $34260449$ and $34260493$ both lie in $(Y,Q]$, so that
$2\cdot34260449\in\vB$ and $2\cdot34260449\cdot34260493\in\vR$, each with small-prime part
$2\le\sqrt Q$; hence $B_{2},R_{2}>0$ and \eqref{eq:X} carries positive weight there. What does hold at
every size we can evaluate is that $Y=L^{6}$ is unusable, so any evaluation of \eqref{eq:X} at such $N$
takes some $Y$ far below $L^{6}$, where the cofactor has only a handful of prime slots.
\end{enumerate}

Neither feature is intrinsic. In Lemma~\ref{lem:rough} the step $|e_{d}|\le\max(t_{d},(Ax)_{d})$ discards
the cancellation between $t_{d}$ and $(Ax)_{d}$ on the rows with a large prime factor; at the parameters
we have examined the resulting bound exceeds the true $\ell^{2}$ norm of $e|_{\Pp(d)>Y}$ by a large
factor, which we have not measured systematically and do not quantify here. Exploiting that cancellation
is the natural route to a better constant and a smaller threshold, and we have not attempted it.

\section{Formalization}\label{sec:lean}

A Lean~4 development accompanies this paper (file \texttt{KappaZero.lean}, toolchain
\texttt{leanprover/\allowbreak lean4:v4.33.1}, mathlib \texttt{v4.33.1}; see the data availability
statement). It formalizes, with no \texttt{sorry} and with no axiom beyond mathlib's \texttt{propext},
\texttt{Classical.\allowbreak choice} and \texttt{Quot.\allowbreak sound}, a selection of the steps of the
proof that are finite and exact --- a selection, not all of them: part~(i) of Lemma~\ref{lem:elem} and
the upper bound of its part~(ii), namely that an $n\le N<(Q+1)^{2}$ has at most one prime
factor above $Q$, that this prime divides $n$ only once, and that $\lfloor N/q_{1}\rfloor\le Q$ --- the
Bertrand lower bound $Q/2-1\le\thr$ and the logarithmic bounds of that lemma are not in the file; the
exactness $d\mid n\Leftrightarrow d\mid a_{Y}(n)$ on $Y$-smooth rows together with the row identity
\eqref{eq:Eid} of Lemma~\ref{lem:E}, dead classes included; the large-prime removal in the $M$--$\pi$ form
it takes at
$d=1$, that is Lemma~\ref{lem:identity}; the Bonferroni identity used in Lemma~\ref{lem:phi} together with
the exact size of its truncation error; the parity split and the ratio bound \eqref{eq:par}; and the
containment $Ax\in K_{Q}$ for $x\ge0$.

Several of these are \emph{generic conditional} lemmas rather than statements about the band of this
paper. The row identity is proved for an arbitrary finite set of squarefree integers; the parity split
and its ratio bound take the two class-count identities as hypotheses, for a supplied inventory of
classes, and derive the algebra from them; the cone lemma assumes a nonnegative coefficient vector.
That is the stronger form to prove, and it is deliberate --- but it also means the development contains
no formal definition of $\kappa(N)$ for the band \eqref{eq:band} and no theorem about it:
\emph{instantiating these lemmas at the paper's actual arithmetic model --- constructing that band
inside Lean, discharging the count identities for it, and assembling the instances --- is not part of
the formal development}.

\emph{The analytic inputs and estimates are not formalized and are not claimed to be}: Mertens' three
theorems, the effective prime number theorem (P), the effective M\"obius bound (L), Rankin's trick, the
sieve, friable and rough-M\"obius estimates of Lemmas~\ref{lem:rank}, \ref{lem:phi}, \ref{lem:U},
\ref{lem:psi} and \ref{lem:roughmob}, the normalisation Proposition~\ref{prop:norm}, and therefore
Theorems~\ref{thm:reduction}, \ref{thm:signed}, \ref{thm:PC} and \ref{thm:main} themselves. The
development certifies the combinatorial skeleton of the proof, not the proof.

\section{A computational remark}\label{sec:numerics}

Because $K_{Q}$ is finitely generated, $\dist_{2}(c_{\vR},K_{Q})$ is the value of a nonnegative
least-squares problem whose data are integral, and its square is rational. Let $\pi$ be the projection
of $c_{\vR}$ on $K_{Q}$, which exists and is unique because $K_{Q}$ is a closed convex cone. If
$\pi=0$ then $\dist_{2}(c_{\vR},K_{Q})^{2}=\|c_{\vR}\|_{2}^{2}$ is an integer. Otherwise write
$\pi=\sum_{\ell}y_{\ell}A_{\ell}$ with $y\ge0$, which is possible because $\pi\in K_{Q}$, and among
all such representations take one whose support $S=\{\ell:y_{\ell}>0\}$ has fewest elements; there are
finitely many generators, so such a representation exists, and $S\neq\emptyset$. Its columns are then
linearly independent. For suppose $\sum_{\ell\in S}v_{\ell}A_{\ell}=0$ with $v\neq0$ supported on $S$;
replacing $v$ by $-v$ if necessary, some $v_{\ell}>0$. Then $y-tv$ represents $\pi$ for every $t$, and
it stays nonnegative for $0\le t\le t^{*}:=\min\{y_{\ell}/v_{\ell}:v_{\ell}>0\}>0$, since the
coordinates with $v_{\ell}\le0$ only increase. At $t=t^{*}$ a further coordinate vanishes, giving a
nonnegative representation of $\pi$ with strictly smaller support and contradicting the choice of $S$.
Let $B$ collect those independent columns and $y_{S}$ the corresponding coefficients. Because
$y_{\ell}>0$ on $S$, the optimality conditions for the projection give
$B^{\mathsf T}(c_{\vR}-\pi)=0$, that is the normal equations
$B^{\mathsf T}B\,y_{S}=B^{\mathsf T}c_{\vR}$ with $B^{\mathsf T}B$ invertible, so
$y_{S}=(B^{\mathsf T}B)^{-1}B^{\mathsf T}c_{\vR}$ and $\pi=By_{S}$ are rational because $B$ and
$c_{\vR}$ are. Hence $\dist_{2}(c_{\vR},K_{Q})^{2}=\|c_{\vR}-\pi\|_{2}^{2}$ is rational, and so is
$\kappa(N)^{2}$ for every $N$ for which $\kappa(N)$ is defined, and both endpoints of an
interval for $\kappa(N)$ can in principle be certified in exact arithmetic. That has been carried out:
both endpoints are certified at every size reached, $N=2^{16},2^{18},\dots,2^{28}$, and at $N=2^{30}$
the lower endpoint is. This section states the inequality the certificates verify and the intervals;
the verification itself is Appendix~\ref{app:verify}.

A pair of certificates suffices. Write $P=c_{\vB}-c_{\vR}$, so that $\kappa(N)=\dist_{2}(c_{\vR},K_{Q})/\|P\|_{2}$
by \eqref{eq:Rfree}. Let $q\ge0$ be an integral vector and $D\ge1$ an integer (a rational primal $q/D$),
and let $\lambda$ be an integral dual with $\theta_{\lambda}(\ell)\ge0$ for \emph{every} $\ell\in\vB$ and
$\langle c_{\vR},\lambda\rangle<0$. Then
\begin{equation}\label{eq:cert}
  \frac{\langle c_{\vR},\lambda\rangle^{2}}{\|P\|_{2}^{2}\,\|\lambda\|_{2}^{2}}
  \ \le\ \kappa(N)^{2}\ \le\ \frac{\|D\,c_{\vR}-Aq\|_{2}^{2}}{D^{2}\,\|P\|_{2}^{2}} ,
\end{equation}
and both bounds are rational and checkable in integer arithmetic from the two stored vectors. The left
inequality is Proposition~\ref{prop:dual} applied to $\lambda/\|\lambda\|_{2}$; the right is
$\dist_{2}(c_{\vR},K_{Q})\le\|c_{\vR}-Aq/D\|_{2}$. The dual feasibility has to be verified on all of
$\vB$, not only on the columns a solver happened to touch; that check is a single pass over the band,
and it is the pass that was run at every size below.

\paragraph{How the two vectors are obtained.} A solver returns a floating-point pair, and neither half
of it is a certificate until it has been rounded to integers and re-checked. The primal is taken to be
$D\,y$ rounded componentwise to the nearest integer and then clamped at $0$, for a per-cell integer
$D$; the clamp keeps $q\ge0$, and the result leaves the right-hand side of \eqref{eq:cert} an exact
rational. At the five smaller two-sided sizes the dual is built from the residual $r=c_{\vR}-Ay$ of the solver's primal, for which $-r$ is dual feasible at the optimum:
$\lambda$ is $-r$ scaled by $K/\max_{d}|r_{d}|$, for a per-cell integer $K$, and rounded componentwise to the nearest integer, so
that $\max_{d}|\lambda_{d}|\le K$. At the two largest no residual is formed. There the column
generation stores a floating dual $w$ with $\langle c_{\vR},w\rangle>0$, already repaired row-wise
before it was saved --- each column whose divisor sum exceeds a threshold is lowered through its
divisor row of least $t_{d}$, never through the row $d=1$, by the largest amount the columns assigned
to that row require, and the threshold kept is the one giving the best floating bound --- and
$\lambda$ is that $-w$, scaled and rounded in the same way. The stored primal is taken from the round
achieving the best upper bound and the stored dual from the round achieving the best lower bound, so
the two halves need not come from the same round, let alone the same vector. The lower-endpoint
certificate at $2^{30}$ is built differently again and
is described in Appendix~\ref{app:verify}. Rounding can make
$\theta_{\lambda}$ negative on a few columns. It is
repaired at the row $d=1$ by a scalar amount $\rho$, which the residual vector $r$ above does not
name: since $1\mid\ell$ for every $\ell$, replacing $\lambda_{1}$ by
$\lambda_{1}+\rho$ adds $\rho$ to $\theta_{\lambda}(\ell)$ at every column at once, so
$\rho=\max(0,-\min_{\ell\in\vB}\theta_{\lambda}(\ell))$ restores feasibility everywhere in one step. The
repair costs something when it does anything at all: since $c_{\vR}(1)=R$ it moves
$\langle c_{\vR},\lambda\rangle$ towards $0$ by $\rho R$, and it changes $\|\lambda\|_{2}^{2}$ by
$2\rho\lambda_{1}+\rho^{2}$, a strict increase exactly when $\rho>0$ \emph{and} the pre-repair $\lambda_{1}$ exceeds $-\rho/2$. Both
conditions matter: the second, which for the stored coordinate $\lambda_{1}+\rho$ reads $\lambda_{1}+\rho>\rho/2$, holds in every certificate stored here, at all
eight sizes, but the first can fail --- a rounded dual that is already feasible has $\rho=0$ and a vacuous
repair, as at $e=20$, where $\min_{\ell\in\vB}\theta_{\lambda}(\ell)$ is already positive. So the
certified lower endpoint is at most the solver's recorded floating optimum --- a nine-decimal ledger
value at the two smallest sizes and, at the next three, the floating optimum stored in the same file
as the dual; those five are the sizes at which the diagnostic clause of Appendix~\ref{app:verify}
compares a stored float against the interval, and at the two largest it compares none --- and
below it where $\rho>0$: an
observation about these certificates, not a property of the repair. The repair needs no new
solve, and the vector stored in the certificate is the repaired one, already feasible as stored.

\paragraph{The certified intervals.} For the seven two-sided cells listed below, both bounds of \eqref{eq:cert} were evaluated in exact integer
arithmetic; at $2^{30}$ only a certified lower bound is available (described below). The endpoints below are the square roots of the two rational bounds that \eqref{eq:cert} gives for $\kappa^{2}$,
rounded outward at twelve decimals; the last column repeats them rounded outward at seven.
\begin{center}
\begin{tabular}{llll}
$e$ & certified $\kappa(2^{e})\ge$ & certified $\kappa(2^{e})\le$ & outward at $7$ decimals \\[2pt]
$16$ & $0.172389351665$ & $0.172389351666$ & $0.1723893<\kappa<0.1723894$ \\
$18$ & $0.140860453831$ & $0.140860453832$ & $0.1408604<\kappa<0.1408605$ \\
$20$ & $0.105542708383$ & $0.105542708385$ & $0.1055427<\kappa<0.1055428$ \\
$22$ & $0.076331197680$ & $0.076331197683$ & $0.0763311<\kappa<0.0763312$ \\
$24$ & $0.052762952627$ & $0.052762952635$ & $0.0527629<\kappa<0.0527630$ \\
$26$ & $0.034622922800$ & $0.034623479647$ & $0.0346229<\kappa<0.0346235$ \\
$28$ & $0.024526262835$ & $0.024538226032$ & $0.0245262<\kappa<0.0245383$ \\
\end{tabular}
\end{center}
The widths are $1.0,1.0,2.0,3.0,8.0$ in units of $10^{-12}$ at $e=16,\dots,24$, then at most
$5.6\cdot10^{-7}$ at $e=26$ and at most $1.2\cdot10^{-5}$ at $e=28$ --- the last two rounded upward, so
the width is below $10^{-6}$ at $2^{26}$ and below $10^{-4}$ at $2^{28}$. How the pair was produced
differs with the size, and the distinction matters only for how
tight the interval is, not for whether it is certified. At $e=16,18,20$ the primal is a dense
nonnegative least squares over the \emph{complete} set of distinct even-$\omega$ columns --- $4\,414$,
$18\,685$ and $76\,946$ of them --- so no column was excluded. At $e=22,\dots,28$ the full problem is
out of reach and column generation was used: a restricted master over a working set, priced by one sweep
over the band per round. Restricting the column set can only increase the distance, so a restricted
optimum, once rounded to an integral $(q,D)$, is still a valid upper bound. At $e=26$ and $28$ the
primal and the dual come from \emph{different} runs: a run stopped on its wall-clock budget leaves a
good primal and a dual weaker than an older converged run's, and \eqref{eq:cert} lets the two halves be
taken from different vectors, each verified on its own. That pairing is what the stored certificates
record, and it is frozen: the numbers above are not to be refreshed from a later re-solve without being
re-certified and re-stated.

\paragraph{The verifier.} The stored certificates are verified in exact integer arithmetic by the
packet's verifier: every decision the check makes is an exact integer or rational comparison --- it
re-derives the cell from the exponent alone, rebuilds the band from scratch, sweeps the dual
feasibility condition over \emph{every} even-$\omega$ column of the whole band, and recomputes both
rational bounds of \eqref{eq:cert}. Appendix~\ref{app:verify} gives those checks in full, together
with the overflow guards, the format and index validation, the one floating-point step and the
measured cost; the packet named in the data availability statement carries the verifier, its run
commands and the frozen timing record, and its README says which outputs belong to which run.

\paragraph{$e=30$: the lower endpoint only.} At $N=2^{30}$ the band has $|\vB|=86\,592\,333$ and
$|\vR|=89\,520\,933$, and only the left half of \eqref{eq:cert} has been evaluated. Writing
\begin{equation}\label{eq:JHD}
  J=-\langle c_{\vR},\lambda\rangle,\qquad H=\|\lambda\|_{2}^{2},\qquad D_{p}=\|P\|_{2}^{2},
\end{equation}
so that the left half of \eqref{eq:cert} reads $\kappa(N)^{2}\ge J^{2}/(H D_{p})$ when $J>0$, the stored
run records the four integers $J$, $H$, $D_{p}$ and $c_{6}$ with
\begin{equation}\label{eq:certint}
  c_{6}^{2}\,H\,D_{p}\ <\ J^{2}\cdot10^{12} ,
\end{equation}
which is that half cleared of denominators and certifies $\kappa(N)>c_{6}/10^{6}$. Here
$J=4401612119558212802>0$, $H=2015370564893177546709684150$, $D_{p}=30382600959846$ and
$c_{6}=17787$, the largest six-decimal value for which \eqref{eq:certint} holds, so
\[
  \kappa(2^{30})\ >\ 0.017787 .
\]
We stress that \eqref{eq:certint} is an arithmetic consequence of feasibility and not a substitute for
it: four integers satisfying an inequality prove nothing on their own. What makes the quantities
\eqref{eq:JHD} the right ones is a stored integral $\lambda$ together with the sweep establishing
$\theta_{\lambda}\ge0$ on every even-$\omega$ column of the whole band, exactly as at the smaller
sizes. The packet ships that $\lambda$ --- the repaired integral vector itself, not the solver's
floating separator --- with the four integers beside it.
No integral primal was
produced at $2^{30}$, so no upper endpoint is claimed there; the restricted master reports
$\kappa(2^{30})\approx0.01779$, which is a solver value and not certified.

\paragraph{The construction's own bounds.} The construction \eqref{eq:X}, evaluated at the same sizes with
$Y$ chosen freely rather than as $L^{6}$ --- which exceeds $Q$ there, see \S\ref{sec:effectivity} ---
gives $\kappa\le0.6982,\,0.6130,\,0.6088,\,0.5841$ at $e=16,18,20,22$, attained at $Y=3,3,5,5$. These
are certified too, and by the cheapest route available: they come from the integer class counts $R_{a}$,
$B_{a}$ through \eqref{eq:X}, so each \emph{squared} bound is an exact rational --- the bounds
themselves are square roots of those rationals and are in fact irrational at all four sizes --- and the
four decimals printed are each the exact square root rounded upward. They exceed the certified optimum by factors between four and eight, and over
the four sizes tested that factor grows with $N$ --- an observation about this range, not something
\S\ref{sec:effectivity} predicts.

One feature of these sizes is worth naming before a reader notices it. Proposition~\ref{prop:norm}
asserts $R\ge0.09N$ and $N/R\to\pi^{2}/(3(1-\log2))=10.7213228\ldots$, both for $L$ large; at every size
computed here the band sits on the other side of those two numbers. The ratio $R/N$ increases with $N$
over the eight sizes, from $0.0768$ at $2^{16}$ to $0.0834$ at $2^{30}$ --- below $0.09$ throughout, and
below the limiting $0.0932720\ldots$ --- so that $N/R$ decreases from $13.03$ to $11.99$, approaching
$10.7213\ldots$ from above and not yet near it. There is no conflict: both
statements of Proposition~\ref{prop:norm} are asymptotic, and a quantity converging to a limit from
above may lie on either side of a weaker constant at any finite point.


The code, the stored primal and dual vectors, and the run commands are in the repository named in the data
availability statement. None of it bears on Theorem~\ref{thm:main} (\S\ref{par:scope}): every size reached
here lies far below the range in which the hypotheses of \S\ref{sec:counting} take hold, so none of these
computations can support or contradict it, and the numbers are reported only because $\kappa(N)$ is a
finite conic distance whose size a reader may wish to know.

\appendix

\section{Verification of the certificates}\label{app:verify}

This appendix records what the check of \S\ref{sec:numerics} does: the cell binding, the format and
index validation, the cache handling, the overflow guards, the one floating-point step and its check,
the measured cost, and the separate construction of the lower-endpoint certificate at the largest
size. Nothing here is needed to read the theorems or \eqref{eq:cert}; it is the account a reader
reproducing the numbers needs.

\paragraph{The verifier.} Each certificate is stored as the two integral vectors together with the
integers they need, and is re-checked by \texttt{pa\_certify.py\ -{}-verify}, which finds and re-checks
every stored certificate; the two largest were produced by \texttt{pa\_certify\_big.py}, which verifies
them too. What the check derives for itself, from the exponent $e$ alone, is the cell: $N=2^{e}$,
$Q=\lfloor\sqrt N\rfloor$, $q_{1}$, $\thr$, and the complete list of squarefree $d\le Q$ from a
M\"obius sieve. It then refuses to proceed unless the cached cell and the stored certificate agree with
that derivation in all five, so a certificate cannot be verified against a cell other than the one it
names. From there the band is rebuilt from scratch and $|\vB|$, $|\vR|$, $c_{\vB}$, $c_{\vR}$ compared
against the cached row counts element for element; $\theta_{\lambda}\ge0$ is swept over \emph{every}
even-$\omega$ column of the whole band; $\langle c_{\vR},\lambda\rangle<0$ and $q\ge0$ are confirmed;
and both rational bounds of \eqref{eq:cert} are recomputed. Every stored index is validated before it
indexes anything and independently of the cached band --- rows in $[1,Q]$, strictly increasing and
squarefree; columns in $(\thr,N]$, distinct, squarefree, $\Pp\le Q$ and of even $\omega$ --- both
integer kernels are guarded by the actual integer sums of the vectors passed to them, and the run
requires the full advertised inventory of seven cells, so a missing certificate fails rather than
quietly shrinking the claim. The stored arrays must also be integer arrays whose values are preserved
by conversion to signed $64$-bit integers; a fractional or value-changing representation is rejected
rather than converted. What is \emph{read} rather than re-derived is the stored pair of integral
vectors itself, which is the object being certified, and the cached per-row counts, which are compared
against the fresh rebuild rather than trusted.

$K$ and
$D$ are chosen per cell so that the integer kernels evaluating \eqref{eq:cert} cannot overflow, and the
totals they produce are checked against exact big-integer arithmetic.

The packet distinguishes a \emph{verification} run, which reads stored certificates and re-checks
them, from a \emph{construction} run, which re-derives and rewrites them; each run announces which it
is, and only the first is a packet check.

It is cheap: one run re-checks all
seven certificates in $73.2$ s at four threads on the authors' machine (times on other hardware
differ) --- $8.6$ s for the five cells $e\le24$
together, $13.2$ s at $e=26$ and $51.4$ s at $e=28$, band rebuild included --- so solving dominates the
cost of these numbers by a wide margin. Re-establishing the $2^{30}$ lower endpoint takes a few minutes
more, the band rebuild dominating there too. The machine is specified with the figures in
the reproduction packet's frozen timing record, which no reproduction step overwrites. The procedure
behind the figures is fixed and the record states it: they are the wall times of one run of the
command, with an empty compiler cache --- a first run, which is what a reader gets, and which pays a
one-time kernel compilation inside the smallest cell --- and the record carries beside them the
totals of the other runs of that command taken on the same machine, cold and with the cache warm, so
that the spread is read off the record and not from any rule for choosing among runs. That spread is
not small: the machine was not idle while any of the runs were taken, carrying unrelated work from
other sessions throughout, and the band rebuild at the largest cell, which sets the total, varies
between runs by more than the compilation the warm runs save. These figures therefore locate the cost
on that machine and are not reproducible to the digit as the arithmetic is.

\paragraph{Two comparisons that are not certificate conditions.} The scripts record two checks beside
\eqref{eq:cert}, reported apart from validity and relied on by nothing. One is a rounding consistency
check: it holds a stored nine-decimal float optimum against the certified interval, and at $e=16,18$
the float lies outside it, a nine-decimal rounding having no reason to lie inside a twelve-decimal
interval. The other is an
internal target: the script behind the two largest certificates asks for a width below $10^{-6}$,
which holds at $2^{26}$ and not at $2^{28}$, where the width is the one printed above --- larger than
the target and inside the $10^{-4}$ claimed there. An unmet target says that the column generation
stopped further from the optimum than was hoped for, not that the certificate is defective: a
certified width is whatever the stored pair gives, and both endpoints are verified in exact arithmetic
either way. The exit codes report certificate validity alone, so neither comparison can be mistaken
for a certificate that failed to verify.

\paragraph{The lower certificate at $2^{30}$.} Its $\lambda$ is \emph{not} built by the normalisation used at the seven two-sided cells, and the
difference matters for anyone comparing the stored vector with the stored separator. Here the solver's
$w$, which the separator holds with $\langle c_{\vR},w\rangle>0$, is rounded componentwise at the
\emph{fixed} scale of \eqref{eq:certint}: $\lambda=\lfloor-10^{12}w\rceil$, the sign being what makes
$\langle c_{\vR},\lambda\rangle<0$ and hence $J>0$. The coordinate at $d=1$ is then \emph{discarded},
set to zero, and rebuilt by the repair, so that $\lambda_{1}$ is exactly the $\rho$ of the repair
described in \S\ref{sec:numerics} and nothing else. Two consequences are worth stating. The scale is not chosen to fill the
integer range: $10^{12}$ is the scale at which the four integers of \eqref{eq:certint} were computed, so
a certificate built this way reproduces the $J$ and $H$ printed above digit for digit instead of
certifying the same bound with different valid integers --- and $\sum_{d}|\lambda_{d}|$ is checked
against the accumulator bound at that scale rather than assumed to fit it. And because
$\max_{d}|\lambda_{d}|$ is here governed by $w$ and the scale rather than set to $K$, it is not the
quantity the seven smaller certificates bound; a reader should expect the stored separator's
$\max_{d}|w_{d}|$ and the certified vector's $\max_{d}|\lambda_{d}|$ to differ by that factor of
$10^{12}$, up to the componentwise rounding and the rebuilt coordinate at $d=1$.

In this constructor --- not in the check, which guards before every kernel it runs --- the sweep that
follows the repair is not preceded by a fresh guard, and does not need one. Write $A$ and $B$ for the
sums of the positive and of the absolute negative coordinates of $\lambda$, so that the quantity
guarded before the first sweep is $\sum_{d}|\lambda_{d}|=A+B$. The first coordinate is zero at that
point, and the repair amount $\rho$ is the negated minimum of $\theta_{\lambda}$ over the even-$\omega$
columns, each value of which is a signed subset sum of the coordinates and so at least $-B$; hence
$0\le\rho\le B$. After the repair every partial sum the kernel forms is $\rho$ plus a signed subset sum
of the remaining coordinates, so it lies between $\rho-B$ and $\rho+A$, hence between $-B$ and $A+B$:
in magnitude it is bounded by the very $A+B$ that was guarded. The vector is in addition required to
satisfy the guard again after the repair, so the certificate that leaves the constructor meets the
blanket condition too.

The check re-establishes the inequality from the stored vector
alone: deriving the cell, rebuilding the band, sweeping every even-$\omega$ column, recomputing $J$,
$H$ and $D_{p}$ from it and confirming that $c_{6}$ is maximal. Its stored arrays are
held to the same format requirement as the others, and no solver is involved: every decision the check
makes is an exact integer or rational comparison, with the integer kernels guarded as above by the
actual integer sums passed to them. Floating-point arithmetic does occur in
the surrounding machinery --- loading the cached arrays and the printed diagnostics, where no verdict depends on it, and
one square root setting a loop bound in the sieve, whose truncation the verifier checks to coincide with the exact integer square root on the whole range where the sieve is used.

\begin{thebibliography}{99}

\bibitem{AlladiJNT}
K.~Alladi,
\emph{Asymptotic estimates of sums involving the Moebius function},
J. Number Theory \textbf{14} (1982), 86--98.

\bibitem{AlladiTAMS}
K.~Alladi,
\emph{Asymptotic estimates of sums involving the Moebius function.\ II},
Trans. Amer. Math. Soc. \textbf{272} (1982), 87--105.

\bibitem{BookerBrowning}
A.~R.~Booker and T.~D.~Browning,
\emph{Square-free values of reducible polynomials},
Discrete Anal. \textbf{2016}, Paper No.~8.

\bibitem{BradyThesis}
Z.~E.~Brady,
\emph{Sieves and iteration rules},
Ph.D.\ thesis, Stanford University, 2017.

\bibitem{BradySDP}
Z.~E.~Brady,
\emph{A semidefinite framework for the sieve},
preprint (2021), \texttt{arXiv:2112.02722}.

\bibitem{BTWirsing}
R.~de~la~Bret\`eche and G.~Tenenbaum,
\emph{Friable averages of oscillating arithmetic functions},
in: Number Theory in Memory of Eduard Wirsing (H.~Maier, J.~Steuding and R.~Steuding, eds.),
Springer, Cham, 2023, pp.~43--71; preprint \texttt{arXiv:2207.04777}, whose listing carries the title
\emph{Friable averages of oscillating multiplicative functions}.

\bibitem{FordMaynard}
K.~Ford and J.~Maynard,
\emph{On the theory of prime producing sieves},
preprint (2024), \texttt{arXiv:2407.14368}.

\bibitem{Franze}
C.~S.~Franze,
\emph{Sifting limits for the $\Lambda^{2}\Lambda^{-}$ sieve},
J. Number Theory \textbf{131} (2011), 1962--1982.

\bibitem{Friedlander}
J.~B.~Friedlander,
\emph{Integers free from large and small primes},
Proc. London Math. Soc. (3) \textbf{33} (1976), 565--576.

\bibitem{OperaDeCribro}
J.~Friedlander and H.~Iwaniec,
\emph{Opera de Cribro},
Amer. Math. Soc. Colloq. Publ.\ \textbf{57}, Amer. Math. Soc., Providence, RI, 2010.

\bibitem{HeathBrown}
D.~R.~Heath-Brown,
\emph{Lectures on sieves},
preprint (2002), \texttt{arXiv:math/0209360}.

\bibitem{HildebrandTenenbaum}
A.~Hildebrand and G.~Tenenbaum,
\emph{Integers without large prime factors},
J. Th\'eor. Nombres Bordeaux \textbf{5} (1993), 411--484.

\bibitem{Polymath}
D.~H.~J.~Polymath,
\emph{Variants of the Selberg sieve, and bounded intervals containing many primes},
Res. Math. Sci. \textbf{1} (2014), Art.~12.

\bibitem{RosserSchoenfeld}
J.~B.~Rosser and L.~Schoenfeld,
\emph{Approximate formulas for some functions of prime numbers},
Illinois J. Math. \textbf{6} (1962), 64--94.

\bibitem{SaiasI}
\'E.~Saias,
\emph{Entiers sans grand ni petit facteur premier.\ I},
Acta Arith. \textbf{61} (1992), 347--374.

\bibitem{SaiasII}
\'E.~Saias,
\emph{Entiers sans grand ni petit facteur premier.\ II},
Acta Arith. \textbf{63} (1993), 287--312.

\bibitem{Soltan}
V.~Soltan,
\emph{Moreau-type characterizations of polar cones},
Linear Algebra Appl. \textbf{567} (2019), 45--62;
preprint \texttt{arXiv:1810.04206}, 2018, whose numbering the theorem cited above follows.

\bibitem{TenenbaumCribler}
G.~Tenenbaum,
\emph{Cribler les entiers sans grand facteur premier},
Philos. Trans. Roy. Soc. London Ser.~A \textbf{345} (1993), 377--384.

\bibitem{TenenbaumBook}
G.~Tenenbaum,
\emph{Introduction to Analytic and Probabilistic Number Theory},
3rd ed., Grad. Stud. Math.\ \textbf{163}, Amer. Math. Soc., Providence, RI, 2015.

\bibitem{Titchmarsh}
E.~C.~Titchmarsh,
\emph{The Theory of the Riemann Zeta-Function},
2nd ed., revised by D.~R.~Heath-Brown, Oxford Univ. Press, Oxford, 1986.

\bibitem{Trudgian}
T.~S.~Trudgian,
\emph{Updating the error term in the prime number theorem},
Ramanujan J. \textbf{39} (2016), 225--234.

\end{thebibliography}

% =====================================================================
\paragraph{Disclosure statement.} No potential competing interest, financial or non-financial, was
reported by the author. No funding was received for this work.
Generative AI tools were used in the derivation, computation, drafting and
auditing of this manuscript; the author takes full responsibility for its
contents. The analytical proofs of the main theorems do not depend on numerical experiments or unverified AI-generated outputs.
The finite-size interval results are supported separately by explicit integer certificates and their documented verification procedures.

\paragraph{Data availability statement.} The audit code, stored outputs, and Lean~4 development accompanying this
manuscript are openly available at
\url{https://github.com/soke91/arithmetic-convolution-fields}, in the
directory \texttt{P7-parity-separation-tends-to-zero/}, at the fixed
release tag \texttt{P7-v1}. No data derived from human or
animal subjects were used.

% =====================================================================
\end{document}
